%% ---------------------------------------------------------------
%%  Z_p-valued residues for regular singular stratified bundles
%%  Working draft.  Compile with pdflatex.
%%
%%  This draft supersedes the version reflected in handoff/zp-chat.md's
%%  "v1"; several of the errors that memo records had already been fixed
%%  here (see \Cref{sec:corrections}).  This revision folds in, from
%%  handoff/zp-residues-v2.tex:
%%    * \S8 rewritten (coordinate independence, Hecke-exponent form of the
%%      boundary twist, replacing the digit-exponent formula, which was
%%      wrong -- see \Cref{sec:corrections}).
%%    * New \S9, p-adic middle convolution, kept explicitly conditional /
%%      conjectural (see status remarks throughout \S9 and
%%      \Cref{sec:corrections}).
%%    * \Cref{thm:localstructure}(2) corrected to "modulo Z" (see
%%      \Cref{sec:corrections}).
%%    * Proof of \Cref{thm:localstructure} (non-single-valued case)
%%      repaired: the k-linear eigenprojectors do not lift to the lattice,
%%      and Hom=0 was conflated with Ext^1=0; replaced by the Z_p/Z-block
%%      decomposition \Cref{lem:Kblocks}.  Statement unchanged.  See
%%      \Cref{sec:corrections}.
%%  \S1--\S7 are otherwise the pre-existing text of this manuscript; where
%%  v2 had simplified a proof that this draft already had in more detail
%%  (e.g.\ \Cref{thm:functoriality}\eqref{it:pullback}), the existing,
%%  more detailed proof was kept.
%% ---------------------------------------------------------------
\documentclass[11pt,reqno]{amsart}

\usepackage{amsmath,amssymb,amsthm,mathrsfs}
\usepackage[T1]{fontenc}
\usepackage{microtype}
\usepackage{stmaryrd}
\usepackage[colorlinks,linkcolor=blue!60!black,citecolor=green!40!black]{hyperref}
\usepackage[capitalise]{cleveref}
\usepackage{xcolor}

%% ---- editorial macros (remove before submission) ---------------
\newcommand{\TODO}[1]{{\color{red}\textsf{[TODO: #1]}}}
\newcommand{\GAP}[1]{{\color{orange}\textsf{[GAP: #1]}}}
\newcommand{\CHECK}[1]{{\color{blue!70!black}\textsf{[check: #1]}}}
\newcommand{\NEEDSLITERATURECHECK}[1]{{\color{green!40!black}\textsf{[NEEDS
LITERATURE CHECK: #1]}}}

%% ---- notation --------------------------------------------------
\newcommand{\Zp}{\mathbb{Z}_p}
\newcommand{\Fp}{\mathbb{F}_p}
\newcommand{\ZZ}{\mathbb{Z}}
\newcommand{\QQ}{\mathbb{Q}}
\newcommand{\NN}{\mathbb{N}}
\newcommand{\GG}{\mathbb{G}}
\newcommand{\PP}{\mathbb{P}}
\newcommand{\AAA}{\mathbb{A}}
\newcommand{\cO}{\mathcal{O}}
\newcommand{\cD}{\mathcal{D}}
\newcommand{\cE}{\mathcal{E}}
\newcommand{\cF}{\mathcal{F}}
\newcommand{\cG}{\mathcal{G}}
\newcommand{\cL}{\mathcal{L}}
\newcommand{\cM}{\mathcal{M}}
\newcommand{\cH}{\mathcal{H}}
\newcommand{\cA}{\mathcal{A}}
\newcommand{\cR}{\mathcal{R}}
\newcommand{\cT}{\mathcal{T}}
\newcommand{\fg}{\mathfrak{g}}
\newcommand{\ft}{\mathfrak{t}}
\newcommand{\Strat}{\mathrm{Strat}}
\newcommand{\rs}{\mathrm{rs}}
\newcommand{\tame}{\mathrm{tame}}
\newcommand{\strat}{\mathrm{strat}}
\newcommand{\Pic}{\mathrm{Pic}}
\newcommand{\Vect}{\mathrm{Vect}}
\newcommand{\Rep}{\mathrm{Rep}}
\newcommand{\Hom}{\mathrm{Hom}}
\newcommand{\End}{\mathrm{End}}
\newcommand{\Ext}{\mathrm{Ext}}
\newcommand{\Spec}{\mathrm{Spec}}
\newcommand{\Frac}{\mathrm{Frac}}
\newcommand{\res}{\mathrm{res}}
\newcommand{\inv}{\mathrm{inv}}
\newcommand{\gr}{\mathrm{gr}}
\newcommand{\Iw}{\mathrm{Iw}}
\newcommand{\dig}{\mathrm{dig}}
\newcommand{\rk}{\mathrm{rk}}
\newcommand{\id}{\mathrm{id}}
\newcommand{\im}{\mathrm{im}}
\newcommand{\tr}{\mathrm{tr}}
\newcommand{\MC}{\mathrm{MC}}
\newcommand{\rig}{\mathrm{rig}}
\newcommand{\Xst}{X_*(T)}
\newcommand{\isoto}{\xrightarrow{\ \sim\ }}
\newcommand{\dR}{\mathrm{dR}}

%% ---- theorem environments --------------------------------------
\theoremstyle{plain}
\newtheorem{theorem}{Theorem}[section]
\newtheorem{proposition}[theorem]{Proposition}
\newtheorem{lemma}[theorem]{Lemma}
\newtheorem{corollary}[theorem]{Corollary}
\newtheorem{conjecture}[theorem]{Conjecture}
\newtheorem{problem}[theorem]{Problem}

\theoremstyle{definition}
\newtheorem{definition}[theorem]{Definition}
\newtheorem{example}[theorem]{Example}
\newtheorem{construction}[theorem]{Construction}

\theoremstyle{remark}
\newtheorem{remark}[theorem]{Remark}
\newtheorem{notation}[theorem]{Notation}

\numberwithin{equation}{section}

%% ---------------------------------------------------------------
\begin{document}

\title[$\Zp$-valued residues]{$\Zp$-valued residues for regular singular
stratified bundles, boundary restriction, and $p$-adic middle convolution}

\author{}
\address{}
\email{}

\date{\today\\ \textsc{Working draft --- not for circulation}}

\begin{abstract}
Let $k$ be an algebraically closed field of characteristic $p>0$, let
$(\bar X,D)$ be a smooth projective $k$-variety with a normal crossings
divisor, and let $U=\bar X\setminus D$.  We attach to every regular
singular stratified bundle on $U$, together with a choice of logarithmic
lattice, a \emph{$\Zp$-valued residue} $r_i$ along each boundary component
$D_i$: a conjugacy class in $\Xst\otimes\Zp$, well defined modulo $\Xst$
independently of the lattice.  The construction uses the fact that the
level-$m$ logarithmic operators $t^{p^m}\partial_t^{[p^m]}$ act with
residues in $\Fp$ which are precisely the $p$-adic digits of $r_i$.

Our main local result (\Cref{thm:localstructure}) is that the category of
regular singular stratified bundles over $k((t))$ is semisimple and
tannakian-equivalent to the representation category of the diagonalisable
group $\Spec k[\Zp/\ZZ]$; the $\Zp$-valued residue, taken \emph{modulo
$\ZZ$}, is a complete invariant of the underlying stratified bundle over
$k((t))$ (\Cref{rem:completemodZ}; an earlier draft omitted the ``modulo
$\ZZ$'', which is corrected in \Cref{sec:corrections}).  In particular the
local monodromy of a regular singular stratified bundle carries no
unipotent part, and any ``limit'' or ``weight'' structure at the boundary
must be of global origin.

Globally we construct a boundary restriction functor
(\Cref{thm:boundary}) decomposing $\bar E_0|_{D_i}$ into $\Zp$-eigen\-com\-po\-nents,
each of which is a regular singular stratified bundle on $D_i$
\emph{twisted by $N_i^{\otimes c}$}, where $N_i$ is the normal bundle and
$c\in\Zp$ the corresponding residue, with the twist governed by an
intrinsic \emph{Hecke exponent} at each level (\Cref{sec:boundary}; an
earlier draft used the wrong, digit-based, exponent here -- see
\Cref{sec:corrections}).  This identifies the $\Zp$-valued residue with a
conformal weight and makes precise an analogy with conformal infinity in
Poincar\'e--Einstein geometry.  The functor is dimension-decreasing and
provides an inductive framework; the argument closes an earlier proof gap
by a coordinate-independence computation that is correct in outline but
still needs a formal write-up (\Cref{sec:boundary}).

Finally, \Cref{sec:MC} sketches a candidate $p$-adic middle convolution
$\MC_c$, $c\in\Zp$, built from Gauss--Manin cohomology.  Because local
monodromy is always semisimple, the classical Jordan-block bookkeeping
disappears from the local invariants entering the construction.  This
part of the paper is exploratory: the rank formula for $\MC_c(\cE)$ is
stated as a \emph{conditional} proposition, made precise on the explicit
hypothesis that $\MC_c$ preserves regular singularity and that a
positive-characteristic Euler-characteristic/middle-extension formula
holds (\Cref{prop:MCrank}); the exponent rules and rigidity invariance of
$\MC_c$ are recorded only as conjectures, verified so far by explicit
low-rank examples (\Cref{sec:MC}).  Whether $\MC_c$ preserves regular
singularity at all is the open \Cref{prob:MCrs}.
\end{abstract}

\maketitle

\setcounter{tocdepth}{1}
\tableofcontents

%% ===============================================================
\section{Introduction}
%% ===============================================================

\subsection{}
In characteristic zero, a flat connection with regular singularities along
a boundary divisor $D$ has a residue whose eigenvalues, read modulo $\ZZ$,
record the local monodromy at infinity; quasi-unipotence of that monodromy
is the fundamental finiteness statement.  Deligne's canonical extension
picks the representative with eigenvalues in $[0,1)$.

In characteristic $p>0$ the correct analogue of a flat connection is a
\emph{stratified bundle} in the sense of Katz \cite{Katz}, i.e.\ an
$F$-divided sheaf $\cE=\bigl((E_n)_{n\ge0},\sigma_n\colon F^*E_{n+1}\isoto
E_n\bigr)$.  Regular singularity was defined and studied by Kindler
\cite{Kindler}, and the associated theory of exponents was developed there
and in Esnault--Kindler \cite{EsnaultKindler}.  The purpose of this paper
is to organise the exponents at \emph{all levels simultaneously} into a
single $p$-adic invariant.

\subsection{The basic mechanism}
Over $\bar X$ the ring of logarithmic differential operators is generated,
in a local coordinate $t$ cutting out $D$, by the operators
\[
   D_a:=t^{a}\partial_t^{[a]},\qquad a\ge0 ,
\]
where $\partial^{[a]}=\partial^a/a!$ is the divided power.  These act on
$t^n$ by $D_a t^n=\binom{n}{a}t^n$; in particular they commute and satisfy
$D_a^{\,p}=D_a$.  Consequently \emph{all} residues of a logarithmic
stratified bundle are semisimple with eigenvalues in $\Fp$
(\Cref{cor:rigidity}).  This rigidity is the source of everything that
follows.

The point is now that the function
\[
   \binom{\cdot}{p^m}\colon \Zp\longrightarrow\Fp
\]
is continuous (it depends only on the argument modulo $p^{m+1}$) and, by
Lucas' theorem, returns the $m$-th $p$-adic digit.  Thus a compatible
family of $\Fp$-valued residues $\bigl(r^{(m)}\bigr)_{m\ge0}$ is nothing
but the digit expansion of a single element
\[
   r=\sum_{m\ge0}r^{(m)}p^{m}\ \in\ \Zp .
\]
The passage from the family of digits to the element of $\Zp$ is not a
formal repackaging: it is exactly the ring structure of $\Zp$ which is
compatible with the geometry (additivity under tensor product,
multiplicativity under Kummer covers, and the residue theorem
$\sum_x r_x=-\deg$).  None of these are visible on the level of digits.

\subsection{Main results}
\begin{itemize}
\item \Cref{thm:welldef}: $r$ is well defined in $\Xst\otimes\Zp$ given a
      level-$0$ logarithmic lattice, and modulo $\Xst$ independently of
      all choices; higher-level lattices do not affect it.
\item \Cref{thm:localstructure} (Local structure theorem): over $k((t))$,
      $\Strat^{\rs}$ is semisimple, equivalent to
      $\Rep\bigl(\Spec k[\Zp/\ZZ]\bigr)$, and $r$ \emph{modulo $\ZZ$} is a
      complete invariant (\Cref{rem:completemodZ}).
\item \Cref{thm:comparison}: $r$ is the limit of the \emph{Iwasawa}
      (flag-wise) Hecke invariants $\nu_n^{\Iw}$; the naive Cartan
      invariants do \emph{not} converge to $r$
      (\Cref{ex:cartanfails}).
\item \Cref{thm:boundary} (Boundary restriction, \Cref{sec:boundary}): a
      canonical $\Zp/\ZZ$-graded decomposition of $\bar E_0|_{D_i}$ into
      $N_i^{\otimes c}$-twisted regular singular stratified bundles on
      $D_i$, with the twist at each level given by an intrinsic Hecke
      exponent (\Cref{cor:heckeexp}).  The argument rests on a
      coordinate-independence statement (\Cref{prop:coordfree}) that is
      correct in outline but not yet fully polished; see the status
      remarks in \Cref{sec:boundary}.
\item \Cref{cor:extension}: $\cE$ extends across $D_i$ as a genuine
      stratified bundle if and only if $r_i\in\Xst$.
\item \Cref{sec:MC} (exploratory): a candidate $p$-adic middle
      convolution $\MC_c$.  \Cref{prop:MCrank} is a conditional
      proposition (rank formula, given \Cref{prob:MCrs} and a
      positive-characteristic Euler-characteristic hypothesis);
      \Cref{prop:MCrules,prop:rigMC} are conjectures, checked only on
      examples; \Cref{prob:MCrs} (does $\MC_c$ preserve regular
      singularity?) is open and is the load-bearing question for all of
      \Cref{sec:MC}.
\end{itemize}

See \Cref{sec:corrections} for a record of corrections to claims made in
earlier drafts of this manuscript.

\subsection{Relation to conformal infinity}
\Cref{thm:boundary} says that the residue eigenvalue $c$ is precisely the
twist by the $c$-th power of the normal bundle.  This is the exact
analogue of a conformal density of weight $c$ on the conformal infinity of
a Poincar\'e--Einstein metric, and explains \emph{why} the invariant must
be $p$-adic: $N_i^{\otimes c}$ only makes sense inside the $F$-divided
formalism for $c\in\Zp$.  See \Cref{sec:conformal}.

\subsection{A negative result}
\Cref{thm:localstructure} shows there is no local unipotent part, hence no
Jacobson--Morozov $\mathfrak{sl}_2$ and no canonical weight filtration
built from local data (\Cref{cor:nounipotent}).  Non-semisimplicity in
$\Strat^{\rs}(U)$ is purely global (\Cref{ex:globalext}).  Any theory of
``limit parabolic geometry at the boundary'' must therefore be built from
$\Ext$-groups over $U$, not from the residue.

%% ===============================================================
\section{Conventions}\label{sec:conventions}
%% ===============================================================

\subsection{}
Throughout $k$ is an algebraically closed field of characteristic $p>0$,
$F$ denotes absolute Frobenius, $\bar X$ is a smooth $k$-variety,
$D=\sum_{i\in I}D_i\subset\bar X$ a strict normal crossings divisor with
smooth irreducible components, and $U=\bar X\setminus D$.  We write
$R=k[[t]]$, $K=k((t))$ for the local situation.

\begin{definition}[Stratified bundles]
A \emph{stratified bundle} on $U$ is a system
$\cE=\bigl((E_n)_{n\ge0},\sigma_n\bigr)$ with $E_n\in\Vect(U)$ and
$\sigma_n\colon F^*E_{n+1}\isoto E_n$.  We set $E:=E_0$.  Writing
$\cD_U^{(n)}\subset\cD_U$ for the subring generated by operators of order
$<p^n$, the isomorphism $E\cong F^{n*}E_n$ endows $E$ with a
$\cD^{(n)}$-horizontal structure, and $\cE$ is the same thing as an
$\cO_U$-coherent $\cD_U$-module which is $\cO$-locally free.
The category $\Strat(U)$ is tannakian over $k$; a fibre functor at a
rational point gives $\pi_1^{\strat}(U)$.
\end{definition}

\begin{definition}[Logarithmic lattices and levels]\label{def:lattice}
A \emph{lattice system} for $\cE$ is a family $\bar E_n\in\Vect(\bar X)$
with $\bar E_n|_U=E_n$.  The maps $\sigma_n$ then induce
$\bar\sigma_n\colon F^*\bar E_{n+1}\to\bar E_n$, isomorphisms away from
$D$, i.e.\ modifications along $D$.  The system is \emph{logarithmic} if
every $\bar E_n$ is stable under $\cD_{\bar X}(\log D)$.
We call $\cE$ \emph{regular singular} if a logarithmic lattice system
exists; this agrees with \cite{Kindler}.
\end{definition}

\begin{notation}\label{not:Da}
In a local coordinate $t$ with $D=\{t=0\}$ put
\[
   D_a:=t^a\partial_t^{[a]}\ \in\ \cD_{\bar X}(\log D),\qquad a\ge0,
   \qquad\qquad \delta^{(m)}:=D_{p^m}.
\]
So $\delta^{(0)}=D_1=t\partial_t$ is the usual Euler operator, and
$\delta^{(m)}$ is the \emph{level-$m$} logarithmic operator: it lies in
$\cD^{(m+1)}$ but not in $\cD^{(m)}$.
For a logarithmic lattice $\bar E_0$ we write
$\theta_a:=\nabla(D_a)\in\End_k(\bar E_0)$ and
$R_a:=\res_D\theta_a\in\End_{\cO_D}(\bar E_0|_D)$, and abbreviate
$R^{(m)}:=R_{p^m}$.
\end{notation}

\begin{definition}[Log parabolic geometry]\label{def:logpar}
Let $G$ be a split reductive group over $\Fp$, $P\subset G$ parabolic,
$T\subset P$ a maximal torus, $W$ the Weyl group, $\fg=\mathrm{Lie}\,G$.
A \emph{log parabolic geometry of type $(G,P)$} on $(\bar X,D)$ is a
$P$-bundle $\cG\to\bar X$ together with a $P$-equivariant isomorphism
$\omega\colon T\cG(-\log\tilde D)\isoto\fg\otimes\cO$ reproducing
fundamental vector fields.  An \emph{infinite-level} log parabolic
geometry is a family $(\cG_n,\omega_n)_{n\ge0}$ such that the adjoint
tractor bundles $\cA_n=\cG_n\times^P\fg$ form a logarithmic lattice system
in the sense of \Cref{def:lattice}.
\end{definition}

\begin{remark}
All residue-theoretic statements below are statements about the filtered
stratified bundle $\cA_\bullet$ alone; non-degeneracy of the Cartan
connection is used only in \Cref{sec:boundary,sec:applications}.  The
reader interested only in $\mathrm{GL}_N$ may take $\cA_\bullet$ to be an
ordinary stratified bundle throughout.
\end{remark}

\subsection{Reduction to the one-variable case}
\begin{lemma}\label{lem:constant}
Let $\bar E_0$ be a logarithmic lattice and $D_i$ connected.  Then the
characteristic polynomial of $R_a$ is constant on $D_i$ and splits over
$\Fp$.
\end{lemma}
\begin{proof}
By \Cref{cor:rigidity} below, at every point the eigenvalues lie in the
discrete set $\Fp$; the multiplicity function is therefore locally
constant, hence constant.
\end{proof}

\begin{corollary}\label{cor:local}
For all local computations we may take a smooth curve germ $C$ transversal
to $D_i$ at a general point and replace $(\bar X,D)$ by
$(\Spec R,\{t=0\})$.
\end{corollary}

%% ===============================================================
\section{Rigidity of logarithmic residues}\label{sec:rigidity}
%% ===============================================================

\begin{proposition}\label{prop:idempotent}
In $\cD_{\bar X}(\log D)$ one has, for all $a,b\ge0$,
\[
   D_a^{\,p}=D_a,\qquad [D_a,D_b]=0 .
\]
\end{proposition}
\begin{proof}
$\cD_{\bar X}(\log D)\subset\End_k(\cO)$ is faithful, so it suffices to
test on the basis $\{t^n\}$ of $k[t,t^{-1}]$ (or on $\{t^n\}_{n\ge0}$ of
$R$, which is again faithful for operators preserving $R$).  From
$\partial^{[a]}t^n=\binom na t^{n-a}$ we get
\[
   D_a\,t^n=\binom{n}{a}\,t^n .
\]
Hence all $D_a$ are simultaneously diagonal in this basis, so they
commute; and the eigenvalues $\binom na\in\Fp$ satisfy $c^p=c$, whence
$D_a^p=D_a$.
\end{proof}

\begin{corollary}\label{cor:rigidity}
Let $\bar E_0$ be a logarithmic lattice for a stratified bundle.  Then the
$k$-linear operators $\theta_a$ satisfy $\theta_a^{\,p}=\theta_a$ and
commute.  In particular the residues $R_a$ are simultaneously
diagonalisable with eigenvalues in $\Fp$.
\end{corollary}

\begin{remark}
\Cref{prop:idempotent} is stronger and more elementary than the usual
argument via Katz's $p$-curvature formula
$\psi(\delta)=\nabla(\delta)^p-\nabla(\delta^{[p]})$, which only handles
$a=1$ directly.  Note that $\theta_a$ is not $R$-linear (only $k$-linear),
so \Cref{cor:rigidity} gives a decomposition of $\bar E_0$ as a
$k$-vector space; the $R$-module structure interacts with it through
\Cref{eq:grading} below.  The coarser grouping of these eigenspaces by
$\Zp/\ZZ$ \emph{does} descend to the lattice and to $K$
(\Cref{lem:Kblocks}); the finer grouping by $c\in\Zp$ does not.
\end{remark}

\begin{lemma}[Leibniz]\label{lem:leibniz}
For $f\in\cO$, $x\in\bar E_0$,
\[
   \theta_a(fx)=\sum_{b+c=a}D_b(f)\cdot\theta_c(x),
\]
and in particular $\theta_a(tx)=t\theta_a(x)+t\,\theta_{a-1}(x)$.
Consequently each $\theta_a$ preserves $t\bar E_0$, so $R_a$ is defined,
and $R_a$ is $\cO_D$-linear.
\end{lemma}
\begin{proof}
The first identity is the divided-power Leibniz rule
$\partial^{[a]}(fx)=\sum_{b+c=a}\partial^{[b]}(f)\partial^{[c]}(x)$
multiplied by $t^a$.  For $f=t$ only $b=0,1$ contribute since
$\partial^{[b]}t=0$ for $b\ge2$.  $\cO_D$-linearity: for $f$ a function in
the directions along $D$ we have $\partial_t^{[b]}f=0$ for $b>0$.
\end{proof}

\begin{lemma}[Lucas relations]\label{lem:lucas}
Write $a=\sum_m a_mp^m$ in base $p$.  Then
\[
   R_a=\prod_{m}\binom{R^{(m)}}{a_m}.
\]
Hence the family $(R_a)_{a\ge0}$ is determined by the level residues
$(R^{(m)})_{m\ge0}$.
\end{lemma}
\begin{proof}
By \Cref{cor:rigidity} we may pass to a simultaneous eigenvector; the
statement becomes Lucas' congruence
$\binom na\equiv\prod_m\binom{n_m}{a_m}\pmod p$.
\end{proof}

%% ===============================================================
\section{The $\Zp$-valued residue}\label{sec:residue}
%% ===============================================================

\subsection{Digits}
\begin{lemma}\label{lem:digitmap}
For each $m\ge0$ the map $x\mapsto\binom{x}{p^m}\bmod p$ extends
uniquely to a continuous map $\Zp\to\Fp$, and this extension sends $x$ to
its $m$-th $p$-adic digit.
\end{lemma}
\begin{proof}
For $n\in\ZZ_{\ge0}$, $\binom{n}{p^m}\bmod p$ depends only on
$n\bmod p^{m+1}$ (Lucas), so the map is uniformly continuous on a dense
subset of $\Zp$ and extends.  Lucas also gives
$\binom{n}{p^m}\equiv n_m$, the $m$-th digit; digits are continuous, so
the identity persists on $\Zp$.
\end{proof}

\begin{definition}[$r$-spectrum]\label{def:rspec}
Let $\cE$ be regular singular on $K$ with logarithmic lattice $\bar E_0$
and let
\[
   \bar E_0/t\bar E_0=\bigoplus_{\chi}\bigl(\bar E_0/t\bar E_0\bigr)_\chi,
   \qquad \chi=(\chi_m)_m\in\Fp^{\NN},
\]
be the simultaneous eigendecomposition of the $R^{(m)}$
(\Cref{cor:rigidity}), only finitely many $\chi$ occurring.  Lifting each
$\chi_m$ to $\{0,\dots,p-1\}$, set
$r_\chi:=\sum_m\chi_mp^m\in\Zp$.  The \emph{$r$-spectrum} is the
multiset
\[
   \mathrm{Spec}_r(\cE,\bar E_0):=
   \bigl\{\,r_\chi\ \text{with multiplicity}\ \dim(\bar E_0/t\bar E_0)_\chi\,\bigr\}
   \subset\Zp .
\]
For a log parabolic geometry of type $(G,P)$ we obtain, after choosing a
maximal torus adapted to the simultaneous diagonalisation, an element
\[
   r\ \in\ \bigl(\Xst\otimes\Zp\bigr)/W
\]
(resp.\ $/W_P$), the \emph{$\Zp$-valued residue}.
\end{definition}

\begin{remark}
The choice of a maximal torus is made once and for all; by
\Cref{cor:rigidity} all levels are diagonalised simultaneously, so there
is a single $W$-ambiguity and not one per level.
\end{remark}

\subsection{Dependence on the lattice}
\begin{theorem}\label{thm:welldef}
Let $\cE$ be regular singular.
\begin{enumerate}
\item The $r$-spectrum depends only on the level-$0$ lattice $\bar E_0$;
      replacing $\bar E_n$ for $n\ge1$ leaves it unchanged.
\item Replacing $\bar E_0$ by a Hecke modification of type $\nu\in\Xst$
      changes $r$ to $r+\nu$.
\item Hence
      $\bar r\in\bigl(\Xst\otimes\Zp/\Xst\bigr)/W$
      is an invariant of $\cE$ alone.
\end{enumerate}
\end{theorem}
\begin{proof}
(1) is immediate from \Cref{def:rspec}, which only refers to $\bar E_0$.
(2) In rank one, $\bar L_0\rightsquigarrow\bar L_0(\nu\cdot D)$ multiplies
the horizontal generator by $t^{-\nu}$, and
$\binom{x+\nu}{p^m}$ is the $m$-th digit of $x+\nu$.  In general reduce to
rank one along the eigendecomposition (or use
\Cref{thm:functoriality}\eqref{it:additive}).

(3) It suffices to show that any two logarithmic lattices $\bar E,\bar E'$
of $\cE$ over $K$ have $r$-spectra that agree modulo $\ZZ$ as multisets.
Replacing the pair by $\bar E+\bar E'\supseteq\bar E,\bar E'\supseteq
\bar E\cap\bar E'$ (both $\theta_a$-stable, hence logarithmic, lattices) we
may assume $\bar E'\subseteq\bar E$.  Then $\bar E/\bar E'$ is a
finite-length $R$-module on which the $\theta_a$ act, semisimply and
commuting by \Cref{cor:rigidity}; every $\theta_a$-simple subquotient is
killed by $t$ (else multiplication by $t$ would be an automorphism of a
finite-length module) and is therefore one-dimensional over $k$, of a
single character $c\in\Zp$.  Choosing a composition series and pulling it
back gives a chain
$\bar E=L_0\supsetneq L_1\supsetneq\dots\supsetneq L_r=\bar E'$ of
$\theta_a$-stable lattices with each $L_i/L_{i+1}$ one such eigenline, of
character $c^{(i)}$.  For $x$ a character-$c^{(i)}$ eigenvector,
\Cref{lem:leibniz} gives
$\theta_b(tx)=t\theta_b(x)+t\theta_{b-1}(x)
=\bigl(\binom{c^{(i)}}{b}+\binom{c^{(i)}}{b-1}\bigr)tx
=\binom{c^{(i)}+1}{b}tx$; comparing $L_i/tL_i$ with $L_{i+1}/tL_{i+1}$,
passing from $L_i$ to $L_{i+1}$ replaces the residue $c^{(i)}$ on that
eigenline by $c^{(i)}+1$ and leaves the other eigenvalues unchanged.
Hence $\mathrm{Spec}_r(\cE,\bar E')$ differs from
$\mathrm{Spec}_r(\cE,\bar E)$ only by integer shifts of individual
eigenvalues, so $\bar r$ modulo $\Xst$ is independent of the lattice; the
$W$-ambiguity is that of \Cref{def:rspec}.
\end{proof}

\begin{remark}
Statement (1) is the precise sense in which ``raising the level shifts the
residue by a Hecke modification'': the level-$m$ residue $R^{(m)}$ is the
$m$-th digit of $r$, so the residues at successive levels are the
successive digits, and only the total $r\in\Zp$ is intrinsic.
\end{remark}

\subsection{Functoriality}
\begin{theorem}\label{thm:functoriality}
Let $\cE,\cE'$ be regular singular with chosen logarithmic lattices.
\begin{enumerate}
\item\label{it:additive} (Tensor)
      $\mathrm{Spec}_r(\cE\otimes\cE')=
      \{\,c+c'\ :\ c\in\mathrm{Spec}_r\cE,\ c'\in\mathrm{Spec}_r\cE'\,\}$,
      sums taken in $\Zp$.
\item (Dual) $\mathrm{Spec}_r(\cE^\vee)=-\mathrm{Spec}_r(\cE)$.
\item (Group homomorphisms) For $\phi\colon G\to G'$, $r'=\mathrm{Lie}(\phi)(r)$;
      compatible with restriction to $P$ and with Levi quotients.
\item\label{it:pullback} (Pullback) If $f\colon(\bar Y,E)\to(\bar X,D)$ is
      such that $f^*D_i=\sum_j e_{ij}E_j$, then
      $r_{E_j}(f^*\cE)=\sum_i e_{ij}\,r_{D_i}(\cE)$.
\end{enumerate}
\end{theorem}
\begin{proof}
\eqref{it:additive}: choose simultaneous eigenvectors $x,x'$ with
$D_a$-eigenvalues $\binom{c}{a},\binom{c'}{a}$ (\Cref{lem:lucas}).  By
\Cref{lem:leibniz} applied to the tensor product,
\[
   \theta_a(x\otimes x')=\sum_{b+c=a}\theta_b(x)\otimes\theta_c(x'),
\]
so the eigenvalue of $D_a$ on $x\otimes x'$ is
$\sum_{b+c=a}\binom{c}{b}\binom{c'}{c}=\binom{c+c'}{a}$ by Vandermonde,
which is valid in $\Zp$ and hence mod $p$.  Taking $a=p^m$ and using
\Cref{lem:digitmap} identifies the residue digits with those of $c+c'$.
(2) is the case $\cE'=\cE^\vee$ of (1) together with
$\mathrm{Spec}_r(\cO)=\{0\}$.
(3) follows from (1),(2) and naturality of $\nabla$.
\eqref{it:pullback}: locally $f^*t=s^{e}\cdot(\text{unit})$, and
$s^{a}\partial_s^{[a]}$ pulls back compatibly; on a horizontal frame
$t^{-A}$ becomes $s^{-eA}$, so the digits are those of $e\,r$.
\end{proof}

\begin{remark}
It is essential that \emph{the eigenvalues add in $\Zp$}, not the digits:
$\dig_m(c)+\dig_m(c')\ne\dig_m(c+c')$ in general because of carries.  This
is the first place where the ring structure of $\Zp$, rather than the
family of digits, is doing real work.
\end{remark}

\begin{corollary}[Tame covers]\label{cor:tame}
Let $[m]$ be a Kummer cover of ramification index $m$ with $p\nmid m$
along $D$.  Then $[m]^*r=m\,r$.
\end{corollary}

%% ===============================================================
\section{The local structure theorem}\label{sec:localstructure}
%% ===============================================================

Throughout this section $R=k[[t]]$, $K=k((t))$.

\begin{definition}
For $c\in\Zp$ with digits $c_n$, let $\cL_c$ be the rank one stratified
bundle on $K$ given by $L_n=K e_n$, $\sigma_n(1\otimes e_{n+1})=t^{-c_n}e_n$.
Then $\bar L_0:=Re_0$ is a logarithmic lattice and
$\mathrm{Spec}_r(\cL_c)=\{c\}$.
\end{definition}

\begin{lemma}\label{lem:mu0}
Let $(\bar E_n,\sigma_n)$ be a lattice system over $R$ in which every
$\sigma_n$ is an isomorphism of lattices.  Then $\cE$ is trivial: $\bar
E_0$ admits a basis of horizontal sections.
\end{lemma}
\begin{proof}
Trivialise $\bar E_n$ and let $g_n\in GL_N(R)$ represent $\sigma_n$.  A
change of basis $v_n$ acts by $g_n\mapsto v_{n+1}^{(p)}g_nv_n^{-1}$, where
$(\cdot)^{(p)}$ is entrywise application of $F$ (a ring endomorphism of
$R$ with $t\mapsto t^p$); trivialising means solving
$v_n=v_{n+1}^{(p)}g_n$.

\emph{Step 1.} Reducing mod $t$, $(\bar g_n)$ is an $F$-divided sheaf on
$\Spec k$; as $k=\bar k$, Lang's theorem makes it trivial (this is the
classical fact that $\Strat(\Spec k)$ is trivial).  After a change of
basis we may assume $g_n\in 1+tM_N(R)$.

\emph{Step 2.} Since $g_{n+i}^{(p^i)}\in1+t^{p^i}M_N(R)$, the product
\[
   v_n:=\lim_{N\to\infty}g_{n+N}^{(p^N)}\cdots g_{n+1}^{(p)}g_n
\]
converges $t$-adically to an element of $1+tM_N(R)\subset GL_N(R)$ and
satisfies $v_n=v_{n+1}^{(p)}g_n$.
\end{proof}

\begin{lemma}\label{lem:res0}
Let $\bar E_0$ be a logarithmic lattice with $R^{(0)}=\res\theta_1=0$.
Then $\bar E_1$ may be replaced so that $\sigma_0(F^*\bar E_1)=\bar E_0$.
\end{lemma}
\begin{proof}
By \Cref{cor:rigidity}, $\theta_1$ is a semisimple $k$-linear operator
with eigenvalues in $\Fp$; write
$\bar E_0=\bigoplus_{c\in\Fp}\bar E_0^{(c)}$.  By \Cref{lem:leibniz},
$\theta_1(tx)=tx+t\theta_1(x)$, hence
\begin{equation}\label{eq:grading}
   t\cdot\bar E_0^{(c)}\subseteq\bar E_0^{(c+1)} :
\end{equation}
$\bar E_0$ is $\Fp$-graded with $\deg t=1$.

The hypothesis says $\theta_1(\bar E_0)\subseteq t\bar E_0$.  If $c\ne0$
and $x\in\bar E_0^{(c)}$ then $cx=\theta_1(x)\in t\bar E_0$, so
$x\in t\bar E_0$.  Therefore $\bar E_0=M+t\bar E_0$ with
$M:=\bar E_0^{(0)}$, and by $t$-adic completeness $\bar E_0=R\cdot M$.

By \eqref{eq:grading}, $t^pM\subseteq M$, so $M$ is a $k[[t^p]]$-module,
and the sum $\sum_{j=0}^{p-1}t^jM$ is direct (the summands lie in distinct
graded pieces) and equals $\bar E_0$; hence $M$ is free of rank $N$ over
$k[[t^p]]$.

Finally $M$ consists of level-$1$ horizontal sections: $\theta_1x=0$ gives
$t\nabla(\partial)x=0$, hence $\nabla(\partial)x=0$ by torsion-freeness,
and $\nabla(\partial^{[j]})=\nabla(\partial)^j/j!=0$ for $1\le j\le p-1$.
Cartier descent applied to $M$ produces $\bar E_1$ with
$\sigma_0(F^*\bar E_1)=R\cdot M=\bar E_0$.
\end{proof}

\begin{lemma}[Level shift]\label{lem:shift}
If $\sigma_0$ is an isomorphism of lattices then $\bar E_1$ is again
$\cD(\log)$-stable, and the level-$m$ residue of
$\cE_1:=(E_{1+i},\sigma_{1+i})$ equals the level-$(m+1)$ residue of $\cE$.
Equivalently $r(\cE)=p\cdot r(\cE_1)$.
\end{lemma}
\begin{proof}
$F^*$ carries $\partial^{[p^m]}$ to $\partial^{[p^{m+1}]}$ up to units,
whence the shift of digits.
\end{proof}

\begin{corollary}\label{cor:allzero}
If all level residues of $\cE$ vanish, i.e.\ $r=0$, then $\cE$ is trivial.
\end{corollary}
\begin{proof}
\Cref{lem:res0} makes $\sigma_0$ a lattice isomorphism; \Cref{lem:shift}
shows $\cE_1$ again has all residues zero; induct, then apply
\Cref{lem:mu0}.
\end{proof}

\begin{lemma}[$\Zp/\ZZ$-blocks over $K$]\label{lem:Kblocks}
Let $\cE$ be regular singular on $K$ with a logarithmic lattice
$\bar E_0$, and $E:=\bar E_0\otimes_RK$.
\begin{enumerate}
\item For every $n$ the commuting semisimple $k$-linear operators
      $\theta_a$ decompose the finite-dimensional space
      $\bar E_0/t^n\bar E_0=\bigoplus_c(\bar E_0/t^n\bar E_0)[c]$ into
      joint eigenspaces, indexed by $c\in\Zp$ with $D_a$ acting by
      $\binom ca$; only finitely many cosets $\bar c\in\Zp/\ZZ$ occur,
      and they are those already occurring in $\bar E_0/t\bar E_0$.
\item Multiplication by $t$ sends $(\bar E_0/t^n\bar E_0)[c]$ into
      $(\bar E_0/t^n\bar E_0)[c+1]$ (Pascal), so the grouping
      $(\bar E_0/t^n\bar E_0)\{\bar c\}
       :=\bigoplus_{c\mapsto\bar c}(\bar E_0/t^n\bar E_0)[c]$ is
      $t$-stable, and
      $\bar E_0/t^n\bar E_0=\bigoplus_{\bar c}(\bar E_0/t^n\bar E_0)\{\bar c\}$
      compatibly in $n$.
\item In the limit,
      $\bar E_0=\bigoplus_{\bar c\in\Zp/\ZZ}\bar E_0\{\bar c\}$ with
      $\bar E_0\{\bar c\}:=\varprojlim_n(\bar E_0/t^n\bar E_0)\{\bar c\}$
      a $\cD_{\bar X}(\log D)$-stable direct summand, hence a sublattice;
      $E\{\bar c\}:=\bar E_0\{\bar c\}\otimes_RK$ is then a
      $\cD_U$-stable $K$-subspace and
      $\cE=\bigoplus_{\bar c\in\Zp/\ZZ}\cE\{\bar c\}$ in
      $\Strat^{\rs}(K)$.  Here $\cE\{\bar c\}$ is intrinsic: it is the
      largest $\cD_U$-stable $K$-subspace of $E$ all of whose
      $\theta_a$-eigen-characters lie in $\bar c$.  Hence the
      decomposition is functorial in $\cE$.
\item The $r$-spectrum of $\cE\{\bar c\}$ lies in the coset $\bar c$;
      after finitely many elementary modifications along its residue
      eigencomponents one obtains a logarithmic lattice of
      $\cE\{\bar c\}$ with a single residue value $c_1\in\bar c$.
\end{enumerate}
The individual joint eigenspaces $E[c]:=\{x\in E:\theta_ax=\binom cax\
\forall a\}$ also make sense, with $t\colon E[c]\isoto E[c+1]$; but as
the $\theta_a$ are $k$-linear, not $R$- or $K$-linear (the Remark after
\Cref{cor:rigidity}), an individual $E[c]$ is neither a $K$-subspace nor
the fibre of a sublattice, and $\bigoplus_{c\in\Zp}E[c]$ is in general a
\emph{proper} $k$-subspace of $E$ (for $\cL_c$ it is
$k[t,t^{-1}]e_0\subsetneq Ke_0$).  Grouping by $\Zp/\ZZ$ -- which
restores stability under $t$ and $t^{-1}$ -- is what yields the genuine
decomposition of (3).
\end{lemma}
\begin{proof}
(1) $\theta_a$ preserves $\bar E_0$ and every $t^j\bar E_0$ (by
\Cref{lem:leibniz}, $\theta_a(t^jx)=\sum_bD_b(t^j)\theta_{a-b}(x)
=\sum_b\binom jb t^j\theta_{a-b}(x)\in t^j\bar E_0$), hence acts on
$\bar E_0/t^n\bar E_0$; there the
$\theta_a$ are commuting and, by \Cref{cor:rigidity}, semisimple with
eigenvalues in $\Fp$, so simultaneously diagonalisable, and by
\Cref{lem:lucas} each joint eigen-character is $a\mapsto\binom ca$ for a
unique $c\in\Zp$ (its digits are the eigenvalues of the
$\theta_{p^m}$), with $D_a=\prod_m\binom{\theta_{p^m}}{a_m}$ acting by
$\binom ca$.  For $0\le j<n$, multiplication by $t^j$ induces an
isomorphism $\bar E_0/t\bar E_0\isoto t^j\bar E_0/t^{j+1}\bar E_0$
shifting $c$ by $j$ (part (2)); so every character of
$\bar E_0/t^n\bar E_0$ has the form $c^{(i)}+j$ with $c^{(i)}$ a
character of $\bar E_0/t\bar E_0$, and lies in the coset $\bar c^{(i)}$.

(2) For a joint eigenvector $x$ of character $c$, \Cref{lem:leibniz}
gives $\theta_a(tx)=t\theta_a(x)+t\theta_{a-1}(x)
=\bigl(\binom ca+\binom c{a-1}\bigr)tx=\binom{c+1}{a}tx$ (Pascal), so
$t$ raises $c$ by $1$; grouping over a coset is thus $t$-stable.  The
transition maps $\bar E_0/t^{n+1}\bar E_0\to\bar E_0/t^n\bar E_0$ are
$\theta_a$-equivariant, hence respect the coset grading.

(3) $\bar E_0=\varprojlim_n\bar E_0/t^n\bar E_0$ ($t$-adic completeness of
a lattice), and a compatible finite direct sum decomposition at each
level passes to the limit, so
$\bar E_0=\bigoplus_{\bar c}\bar E_0\{\bar c\}$.  Each
$\bar E_0\{\bar c\}$ is $\cD_{\bar X}(\log D)$-stable (the $\theta_a$
respect the coset grading at every level) and is an $R$-module direct
summand of the finite free module $\bar E_0$, hence finite free -- a
sublattice.  Then $E\{\bar c\}=\bar E_0\{\bar c\}\otimes_RK$ is a
$K$-subspace with $E=\bigoplus_{\bar c}E\{\bar c\}$, and it is
$\cD_U$-stable: it is $K$-stable, and the operators
$\nabla(\partial_t^{[j]})=t^{-j}\theta_j$ preserve it because $\theta_j$
respects the coset grading (it acts on each character space by a scalar)
and $t^{-j}$ preserves the $K$-subspace.  A $\cD_U$-submodule of $\cE$
that is a $K$-subspace is a sub-stratified bundle, and the $r$-spectrum of
$\cE\{\bar c\}$ lies in $\bar c$ by construction.

\emph{Intrinsic characterisation.}  We use that $\Strat^{\rs}(K)$ is
closed under subquotients (\cite{Kindler}; it is a tannakian subcategory
of $\Strat(K)$), that a nonzero regular singular object has a
nonzero $r$-spectrum ($\bar V\ne0\Rightarrow\bar V/t\bar V\ne0$ by
Nakayama), and that for a surjection $V\twoheadrightarrow W$ of
regular singular objects the $r$-spectrum of $W$ lies (mod $\ZZ$) in that
of $V$ (choose lattices $\bar V\twoheadrightarrow\bar W$; the
$\theta_a$-semisimple $\bar W/t\bar W$ is a $\theta_a$-equivariant
quotient of $\bar V/t\bar V$).  Now let $V\subseteq E$ be any
$\cD_U$-stable $K$-subspace with $r$-spectrum in $\bar c$.  For
$\bar d\ne\bar c$ the composite $V\hookrightarrow E\twoheadrightarrow
E\{\bar d\}$ has image both a subobject of $E\{\bar d\}$ ($r$-spectrum in
$\bar d$) and a quotient of $V$ ($r$-spectrum in $\bar c$); as
$\bar c\cap\bar d=\varnothing$ that image is $0$.  Hence
$V\subseteq\bigcap_{\bar d\ne\bar c}\ker(E\to E\{\bar d\})=E\{\bar c\}$.
This proves maximality and lattice-independence of $E\{\bar c\}$.
Functoriality follows: for $\varphi\colon\cE\to\cE'$ in $\Strat^{\rs}(K)$,
$\varphi(E\{\bar c\})$ is a $\cD_U$-stable $K$-subspace with $r$-spectrum
in $\bar c$, so $\varphi(E\{\bar c\})\subseteq E'\{\bar c\}$.

(4) The $r$-spectrum of $\bar E_0\{\bar c\}$ is the sub-multiset of that
of $\bar E_0$ carried by the characters with $c\mapsto\bar c$, so lies in
$\bar c$; independence modulo $\ZZ$ of the lattice is
\Cref{thm:welldef}(2).  The residue values in the block differ by
integers, so finitely many elementary modifications along the residue
eigencomponents equalise them to a single $c_1\in\bar c$: this is the
case $R=k[[t]]$ of \Cref{lem:normalise} (equivalently \Cref{prop:jump}),
proved in \Cref{sec:boundary} by an elementary Vandermonde computation
that uses neither the present lemma nor \Cref{thm:localstructure}.

For the closing remark: the eigenvector computation of (2) gives
$t\colon E[c]\to E[c+1]$, and applying \Cref{lem:leibniz} to
$t\cdot(t^{-1}y)=y$ for $y\in E[c+1]$, then inducting on $a$ (with
$\theta_0=\id$), gives $t^{-1}\colon E[c+1]\to E[c]$; so
$t\colon E[c]\isoto E[c+1]$, and $\bigoplus_cE[c]$ is a $k$-subspace
stable under $t^{\pm1}$.  For $\cL_c$ it is $k[t,t^{-1}]e_0$, properly
contained in $Ke_0$, so this $k$-subspace is in general not all of $E$
and does not by itself give a $K$-linear decomposition.
\end{proof}

\begin{theorem}[Local structure theorem]\label{thm:localstructure}
Let $\cE$ be a regular singular stratified bundle on $K=k((t))$.  Then
there is an isomorphism
\[
   \cE\ \cong\ \bigoplus_{c\in\Zp}\cL_c^{\oplus n_c}\qquad(\text{finite sum}).
\]
Consequently:
\begin{enumerate}
\item $\Strat^{\rs}(K)$ is a semisimple tannakian category, equivalent to
      $\Rep\bigl(\cT_p\bigr)$ where
      $\cT_p:=\Spec k[\Zp/\ZZ]$ is diagonalisable with character group
      $\Zp/\ZZ$;
\item the $r$-spectrum \emph{modulo $\ZZ$} is a complete invariant
      (\Cref{rem:completemodZ});
\item $\Ext^1_{\Strat^{\rs}(K)}(\cL_c,\cL_{c'})=0$ for all $c,c'$.
\end{enumerate}
\end{theorem}
\begin{proof}
\emph{Single value.}  Suppose the $r$-spectrum is a single value $c$ with
multiplicity $N$.  Then every $R^{(m)}$ is the scalar $\dig_m(c)$, so by
\Cref{thm:functoriality}\eqref{it:additive} the bundle
$\cL_{-c}\otimes\cE$ has $r$-spectrum $\{0\}$ with multiplicity $N$; by
\Cref{cor:allzero} it is trivial, whence $\cE\cong\cL_c^{\oplus N}$.

\emph{General case.}  By \Cref{lem:Kblocks}(3) there is a finite
decomposition $\cE=\bigoplus_{\bar c\in\Zp/\ZZ}\cE\{\bar c\}$ in
$\Strat^{\rs}(K)$.  Fix $\bar c$.  By \Cref{lem:Kblocks}(4) some
logarithmic lattice of $\cE\{\bar c\}$ has a single residue value
$c_1\in\bar c$, so the single-value case applied to $\cE\{\bar c\}$ gives
$\cE\{\bar c\}\cong\cL_{c_1}^{\oplus n_{\bar c}}$ with
$n_{\bar c}=\dim_KE\{\bar c\}$.  Summing over $\bar c$ gives
$\cE\cong\bigoplus_{\bar c}\cL_{c(\bar c)}^{\oplus n_{\bar c}}$, an
isomorphism of the asserted form $\bigoplus_{c\in\Zp}\cL_c^{\oplus n_c}$.

This establishes the displayed isomorphism; \emph{every} object of
$\Strat^{\rs}(K)$ is a finite direct sum of the $\cL_c$.  Statements (1)
and (3) are now formal consequences, as follows.

\emph{Morphisms between the $\cL_c$.}  For $c\equiv c'\pmod\ZZ$ we have
$\cL_c\cong\cL_{c'}$ (\Cref{rem:completemodZ}) and
$\End(\cL_c)=k$: a $\cD_U$-linear endomorphism of the rank-one
$K$-vector space $\cL_c$ is multiplication by a horizontal function,
i.e.\ by a scalar.  For $c\not\equiv c'\pmod\ZZ$, $\Hom(\cL_c,\cL_{c'})=0$:
a nonzero map between rank-one objects over the field $K$ is an
isomorphism, forcing $\cL_c\cong\cL_{c'}$ and hence $[c]=[c']$ in
$\Zp/\ZZ$ by \Cref{thm:welldef}(3) (the residue mod $\ZZ$ is an
isomorphism invariant), a contradiction.  So the $\cL_{\bar c}$, one for each coset
$\bar c\in\Zp/\ZZ$, are pairwise non-isomorphic objects with
endomorphism ring $k$ and no nonzero maps between distinct ones.

\emph{Statement (1).}  Combining the two preceding paragraphs, the
$k$-linear functor
$\omega\colon\cE\mapsto\bigl(\Hom(\cL_c,\cE)\bigr)_{\bar c\in\Zp/\ZZ}$,
to finitely supported $\Zp/\ZZ$-graded finite-dimensional $k$-vector
spaces, is fully faithful (both sides compute
$\prod_{\bar c}\mathrm{Mat}_{m_{\bar c}\times n_{\bar c}}(k)$ on
$\cE=\bigoplus\cL_c^{n_c}$, $\cE'=\bigoplus\cL_c^{m_c}$) and essentially
surjective; hence an equivalence of additive categories.  Under it
$\Strat^{\rs}(K)$ inherits from the target the structure of a semisimple
abelian $k$-linear category, and the tensor product corresponds to the
convolution of $\Zp/\ZZ$-gradings, so
$\Strat^{\rs}(K)\simeq\Rep(\cT_p)$ with $\cT_p:=\Spec k[\Zp/\ZZ]$: the
representations of the diagonalisable group scheme $\Spec k[A]$ of an
abelian group $A$ are exactly the $A$-graded $k$-vector spaces, the
finite-dimensional ones having finite weight support
\NEEDSLITERATURECHECK{standard; e.g.\ SGA3 Exp.\ I or Waterhouse,
\emph{Introduction to Affine Group Schemes}, ch.\ 2 -- add precise
reference}.  Since $\Zp/\ZZ$ has no $p$-torsion (its torsion subgroup
$\ZZ_{(p)}/\ZZ$ has order prime to $p$), $\cT_p$ is a genuine
diagonalisable group and $\Rep(\cT_p)$ -- graded vector spaces -- is
semisimple.

\emph{Statement (3).}  In $\Rep(\cT_p)$ -- equivalently in
$\Zp/\ZZ$-graded vector spaces -- all $\Ext^1$ vanish; transporting along
the equivalence $\omega$, $\Ext^1_{\Strat^{\rs}(K)}(\cL_c,\cL_{c'})=0$ for
\emph{all} $c,c'$.  Concretely, an extension $0\to\cL_{c'}\to M\to\cL_c\to0$
has $M\cong\bigoplus_d\cL_d^{\oplus m_d}$ by the first part of the proof;
the inclusion $\cL_{c'}\hookrightarrow M$ is a nonzero element of
$\Hom(\cL_{c'},M)$, which by the $\Hom$-computation is the multiplicity
space of the $\bar c'$-isotypic summand $M_{\bar c'}\cong\cL_{c'}^{\oplus m_{\bar c'}}$
of $M$; its image is a line, a direct summand of $M_{\bar c'}$ and hence
of $M$, and the complementary summand maps isomorphically to $\cL_c$
under $M\to\cL_c$, giving a section.  \emph{This uses the direct-sum
decomposition of the extension object, not $\Hom(\cL_c,\cL_{c'})=0$, which
does not by itself force splitting; see \Cref{sec:corrections}.}

\emph{Statement (2)} is \Cref{rem:completemodZ}.
\end{proof}

\begin{remark}[Isomorphisms among the $\cL_c$, and consistency with
\Cref{prop:Gm}]\label{rem:completemodZ}
For $n\in\ZZ=\Xst$ (rank one) we have $\cO\cong\cL_n$: equip $\cO$ with the
Hecke-modified lattice, $t^{-n}$ times the standard generator at level
$0$, leaving the higher lattices unchanged; by \Cref{thm:welldef}(2) this
lattice has $r$-spectrum $\{n\}$, and since $\cO$ has rank one this is the
single-eigenvalue case handled in the first paragraph of the proof above,
so $\cO\cong\cL_n$.  Tensoring by $\cL_c$ and applying the same
single-eigenvalue case to $\cL_c\otimes\cL_n$ (whose $r$-spectrum is
$\{c+n\}$ by \Cref{thm:functoriality}\eqref{it:additive}) gives
\[
   \cL_c\ \cong\ \cL_c\otimes\cO\ \cong\ \cL_c\otimes\cL_n\ \cong\ \cL_{c+n},
   \qquad n\in\ZZ .
\]
So the $\cL_c$, $c\in\Zp$, are pairwise isomorphic exactly along
$\ZZ$-cosets.  This is why item (2) above must read ``modulo $\ZZ$'' and
not literally $\Zp$: it is the same phenomenon, one dimension down, as
$\Pic^{\strat}(\GG_m)\cong\Zp/\ZZ$ in \Cref{prop:Gm} -- an integral shift
of the level-$0$ lattice generator, over a field where $t$ is invertible,
is already an isomorphism of the underlying object, not merely a change
of some auxiliary integral structure.  It is also forced by item (1) of
the theorem: representations of the diagonalisable group
$\cT_p=\Spec k[\Zp/\ZZ]$ are graded by characters valued in $\Zp/\ZZ$, not
$\Zp$, so a complete invariant for $\Strat^{\rs}(K)\simeq\Rep(\cT_p)$ can
only be the $r$-spectrum read modulo $\ZZ$; item (2) as originally stated
(without ``modulo $\ZZ$'') was inconsistent with item (1) of the very same
theorem.  See \Cref{sec:corrections} for the correction history.

None of this affects the decomposition
$\cE\cong\bigoplus_{c\in\Zp}\cL_c^{\oplus n_c}$ itself, which remains a
valid (if redundant, since $c$ and $c+n$ give the same summand up to
isomorphism) way of exhibiting $\cE$, nor does it affect
\Cref{cor:extension}, which concerns a fixed lattice attaining residue
exactly $0$, not the isomorphism class of $\cE$ over $K$.
\end{remark}

\begin{corollary}[No local unipotent part]\label{cor:nounipotent}
The local monodromy of a regular singular stratified bundle is semisimple.
There is no nilpotent element of $\fg$ attached to $r$, hence no
Jacobson--Morozov $\mathfrak{sl}_2$ and no canonical weight filtration
constructed from local data.
\end{corollary}
\begin{proof}
The residues $R_a$ are simultaneously diagonalisable with eigenvalues in
$\Fp$ (\Cref{cor:rigidity}), so the local monodromy is semisimple.  By
\Cref{thm:localstructure} every object of $\Strat^{\rs}(K)$ is
$\bigoplus_c\cL_c^{\oplus n_c}$ and the category is semisimple, so no
local datum records a nilpotent (Jordan-block) structure: the only local
invariant is the $r$-spectrum modulo $\ZZ$, a multiset of semisimple
characters.
\end{proof}

\begin{example}[Non-semisimplicity is global]\label{ex:globalext}
On $\GG_m=\Spec k[t,t^{-1}]$ consider the extension of $\cO$ by $\cO$
given by $f_n=t$, i.e.\ $\sigma_n=\begin{pmatrix}1&t\\0&1\end{pmatrix}$.
Splitting requires $g_n-g_{n+1}^p=t$ with $g_n\in k[t,t^{-1}]$.  If
$d_n:=\deg g_n$ then $d_n=\max(1,pd_{n+1})$; if some $d_{n+1}\ge1$ the
degrees are unbounded going down, and if all $d_n=0$ then $g_n\in k$
contradicts $g_n=t+g_{n+1}^p$.  So the extension is non-split, although
$r=0$.  Over $k[[t]]$, however, $g=\sum_{i\ge0}t^{p^i}$ solves
$g=g^p+t$, so it splits, in accordance with
\Cref{thm:localstructure}.
\end{example}

\begin{remark}
\Cref{thm:localstructure} should be compared with Kindler's theory of
exponents \cite{Kindler}; the overlap, and whether semisimplicity is
stated there, must be determined before circulation.
\TODO{compare with \cite{Kindler}, \cite{EsnaultKindler}}
\end{remark}

%% ===============================================================
\section{Comparison with Hecke invariants}\label{sec:comparison}
%% ===============================================================

\subsection{The horizontal lattices $\Sigma_n$}
Set
\[
   \tau_n:=\sigma_0\circ F^*\sigma_1\circ\cdots\circ F^{(n-1)*}\sigma_{n-1}
   \colon F^{n*}\bar E_n\to E:=\bar E_0\otimes K,
   \qquad \Sigma_n:=\im\tau_n .
\]
Thus $\Sigma_0=\bar E_0$ and $\Sigma_n$ is spanned by level-$n$ horizontal
vectors.

\begin{lemma}\label{lem:sigmainv}
$\inv(\Sigma_n,\Sigma_{n+1})=p^n\mu_n$, where
$\mu_n=\inv\bigl(\bar E_n,\sigma_n(F^*\bar E_{n+1})\bigr)$.
\end{lemma}
\begin{proof}
$F^{n*}$ multiplies valuations by $p^n$.
\end{proof}

\begin{proposition}\label{prop:flagres}
Suppose $\tau_n$ is upper triangular with diagonal entries
$t^{A^{(1)}_n},\dots,t^{A^{(N)}_n}$ with respect to a basis
$e_1,\dots,e_N$ of $\bar E_0$ and some basis of $\bar E_n$.  Then for
$m<n$, $\theta_{p^m}$ preserves the flag $\langle e_1,\dots,e_j\rangle$ and
\[
   \theta_{p^m}(e_j)\equiv\binom{-A^{(j)}_n}{p^m}e_j
   \pmod{\langle e_1,\dots,e_{j-1}\rangle+t\bar E_0}.
\]
\end{proposition}
\begin{proof}
Let $f_i:=\tau_n(\text{$i$-th basis vector})$, so $\partial^{[j]}f_i=0$ for
$0<j<p^n$.  Triangularity gives
$e_j=t^{-A^{(j)}_n}f_j+\sum_{i<j}c_{ij}f_i$ with $c_{ij}\in K$.  For
$m<n$ the divided-power Leibniz rule collapses to
$\partial^{[p^m]}(cf)=\partial^{[p^m]}(c)f$, whence
$\theta_{p^m}(e_j)=\binom{-A^{(j)}_n}{p^m}e_j
+\sum_{i<j}\bigl(\delta^{(m)}c_{ij}\bigr)f_i$.
\end{proof}

\subsection{The Iwasawa invariant}
\Cref{thm:localstructure} guarantees a full flag exists (indeed a
splitting).  Define
\[
   \nu_n^{\Iw}:=\bigl(A^{(1)}_n,\dots,A^{(N)}_n\bigr)\in\Xst ,
\]
the diagonal exponents of $\tau_n$ with respect to a fixed full flag.

\begin{theorem}\label{thm:comparison}
$\nu^{\Iw}$ is strictly additive,
$\nu^{\Iw}_{n+1}=\nu^{\Iw}_n+p^n\mu^{\Iw}_n$; in particular
$\nu^{\Iw}_{n+1}\equiv\nu^{\Iw}_n\pmod{p^n}$, the sequence converges in
$\Xst\otimes\Zp$, and
\[
   -\lim_{n\to\infty}\nu^{\Iw}_n\ =\ r ,
\]
the $r$-spectrum arranged along the flag.
\end{theorem}
\begin{proof}
Additivity is the statement that diagonal entries of a product of
triangular matrices multiply.  Convergence is $p^n\mu_n\to0$.  The
identification with $r$ is \Cref{prop:flagres} combined with
\Cref{lem:digitmap}: $\binom{-A^{(j)}_n}{p^m}\to\binom{r_j}{p^m}=\dig_m(r_j)$.
\end{proof}

\subsection{Two counterexamples}
It is tempting to define $r$ directly from the Cartan (dominant) relative
positions $\mu_n$ in the affine Grassmannian by $r:=-\sum_n\mu_np^n$.
This is wrong.

\begin{example}[Cartan invariants are not additive]\label{ex:cartanfails}
Let $G=\mathrm{GL}_2$, $g_0=\mathrm{diag}(t,t^{-1})$,
$g_1=\mathrm{diag}(t^{-1},t)$, $g_n=1$ for $n\ge2$.  Then
$\mu_0=\mu_1=(1,-1)$ in dominant form, so
$\sum\mu_np^n=(1+p,-1-p)$.  But
\[
   \tau_2\leftrightarrow g_0\cdot g_1^{(p)}=\mathrm{diag}(t^{1-p},t^{p-1}),
   \qquad \nu_2=(p-1,1-p)=\nu_n\ (n\ge2).
\]
The bundle is $\cL_{p-1}\oplus\cL_{1-p}$, so $r=(p-1,1-p)$, matching
$-\lim\nu^{\Iw}_n$ but not $-\sum\mu_np^n$; for $p>3$ the two do not even
agree modulo $p$ up to $W$.  The discrepancy is the positive coroot
correction in the convolution
$\inv(\Sigma_0,\Sigma_2)\le\mu_0+p\mu_1$, which is not divisible by
$p$.
\end{example}

\begin{example}[Convergence of $\nu_n$ does not imply regularity]\label{ex:irregular}
Let $\sigma_n=\begin{pmatrix}1&t^{-1}\\0&1\end{pmatrix}$ for all $n$.
Then $\tau_n=\begin{pmatrix}1&\sum_{i<n}t^{-p^i}\\0&1\end{pmatrix}$ and
$\nu_n=(-p^{n-1},p^{n-1})\to0$.  The bundle is not trivial: splitting
needs $h_{n+1}^p-h_n=t^{-1}$ in $K$, and pole orders satisfy
$d_n=\max(pd_{n+1},1)$, forcing unboundedness.  By
\Cref{thm:localstructure} it is therefore \emph{not} regular singular.
Thus the logarithmic condition in \Cref{def:lattice} cannot be dropped and
$\lim\nu_n$ is not by itself a residue.
\end{example}

%% ===============================================================
\section{Rank one and low-rank examples}\label{sec:examples}
%% ===============================================================

\subsection{$\GG_m$}
\begin{proposition}\label{prop:Gm}
$\Pic^{\strat}(\GG_m)\cong\Zp/\ZZ$, via $\cL\mapsto\bar r_0$.
\end{proposition}
\begin{proof}
Since $\Pic(\GG_m)=0$, write $L_n=\cO e_n$ and
$\sigma_n(1\otimes e_{n+1})=c_nt^{a_n}$, $c_n\in k^\times$.  A change
$e_n'=d_nt^{b_n}e_n$ gives $c'_n=c_nd_{n+1}^pd_n^{-1}$ and
$a'_n=a_n+pb_{n+1}-b_n$.  As $k=\bar k$ the $c_n$ can be normalised to
$1$.  Set $r:=-\sum_na_np^n$.  Then
$r'-r=-\sum_n(pb_{n+1}-b_n)p^n=b_0-\lim_Nb_Np^N=b_0$, giving
well-definedness modulo $\ZZ$ and $r\mapsto r+b_0$ under lattice change.
Surjectivity: read off digits.  Injectivity: if $\sum a_np^n=0$ then
$A_n=\sum_{i<n}a_ip^i$ is divisible by $p^n$, and $b_n:=A_n/p^n\in\ZZ$
satisfies $pb_{n+1}-b_n=a_n$.
\end{proof}

\begin{proposition}[$p$-adic residue theorem]\label{prop:restheorem}
For $\bar X=\PP^1$, $D=\{0,\infty\}$ and a rank one $\cE$ with lattice
$\bar L_0$,
\[
   r_0+r_\infty=-\deg\bar L_0\ \in\ \ZZ\subset\Zp .
\]
\end{proposition}
\begin{proof}
In the coordinate $s=t^{-1}$ the exponents change sign, so
$r_\infty=-r_0$ when $\bar L_0=\cO$.  Twisting by a point changes
$\deg$ by $1$ and one residue by $\mp1$ (\Cref{thm:welldef}(2)).
\end{proof}

\begin{proposition}[Kummer interpretation]\label{prop:kummer}
Under \Cref{prop:Gm}:
\begin{enumerate}
\item $\bar r=0$ iff $\cE$ is trivial;
\item $\bar r$ is torsion, i.e.\ $r\in\ZZ_{(p)}$, iff $\cE$ is a Kummer
      sheaf $\cL_\chi$ for some $\chi\colon\mu_m\to k^\times$,
      $p\nmid m$; indeed
      $\ZZ_{(p)}/\ZZ\cong(\QQ/\ZZ)^{(p')}\cong
      \Hom(\pi_1^{\tame}(\GG_m),k^\times)$;
\item if $r\in\Zp\setminus\ZZ_{(p)}$ then $\cE$ is a $p$-adic limit of
      Kummer sheaves not coming from $\pi_1^{\tame}$.
\end{enumerate}
\end{proposition}

\begin{corollary}[Why canonical extensions are a tame phenomenon]\label{cor:nosection}
The projection $\Zp\to\Zp/\ZZ$ admits no continuous section; hence there
is in general no canonical choice of lattice normalising $r$.  On the
torsion part, however, $[0,1)\cap\ZZ_{(p)}$ is a canonical section.  This
is exactly the range in which Katz's canonical extension \cite{Katz}
exists.
\end{corollary}

\subsection{$\PP^1\setminus\{0,1,\infty\}$}
\begin{proposition}
$\Pic^{\strat}(\PP^1\setminus\{0,1,\infty\})\cong(\Zp/\ZZ)^2$ via
$(\bar r_0,\bar r_1)$, with $\bar r_\infty=-\bar r_0-\bar r_1$.
\end{proposition}

\begin{example}[Irreducible non-tame rank two]\label{ex:GL2}
Let $p\ne2$ and $f\colon\PP^1_z\to\PP^1_t$, $t=z^2$, branched at
$0,\infty$; $f^{-1}(1)=\{\pm1\}$.  Let $\cM$ be rank one regular singular
on $\PP^1_z\setminus\{0,1,-1,\infty\}$ with residues
$\alpha,\beta,\gamma,\delta$.  Then $\cE:=f_*\cM$ is a rank two regular
singular stratified bundle on $U=\PP^1\setminus\{0,1,\infty\}$ with
\[
   r_0=\Bigl\{\tfrac\alpha2,\tfrac{\alpha+1}2\Bigr\},\quad
   r_1=\{\beta,\gamma\},\quad
   r_\infty=\Bigl\{\tfrac\delta2,\tfrac{\delta+1}2\Bigr\},
\]
(note $\tfrac12\in\Zp$ as $p\ne2$), and $\cE$ is irreducible whenever
$\cM\not\cong\iota^*\cM$ for $\iota\colon z\mapsto-z$.  The residue
theorem checks out: $\sum_x\tr r_x=\alpha+\beta+\gamma+\delta+1$, and
$\det f_*\cO=\cO(-1)$ accounts for the $+1$.
Taking $\alpha\notin\ZZ_{(p)}$ produces an irreducible regular singular
stratified bundle which is not tame.
\end{example}

\begin{corollary}
The kernel of $\pi_1^{\strat,\rs}(U)\to\pi_1^{\tame}(U)$ is nontrivial
already for $U=\PP^1\setminus\{0,1,\infty\}$.
\end{corollary}

\begin{problem}[$p$-adic rigidity]\label{prob:rigid}
Is an irreducible rank two regular singular stratified bundle on
$\PP^1\setminus\{0,1,\infty\}$ determined up to isomorphism by its
$r$-spectra $(r_0,r_1,r_\infty)$?  Which triples with
$\sum\tr r_x\in\ZZ$ occur?  Formally Katz's rigidity index
$\mathrm{rig}=(2-2g)N^2-\sum_x(N^2-\dim Z(r_x))=2$ for regular semisimple
$r_x$, and $Z(r_x)$ makes sense as the joint centraliser of the digits.
\end{problem}

\begin{remark}
\Cref{ex:GL2} only realises triples of the special shape above.  A
construction of general triples would presumably be a $p$-adic middle
convolution.  \Cref{thm:localstructure} is encouraging here: since there
are no nontrivial local extensions, the intermediate extension $j_{!*}$
may be unambiguous even without a normalisation of $r$ into $[0,1)$,
which \Cref{cor:nosection} forbids.
\TODO{develop $p$-adic middle convolution}
\end{remark}

%% ===============================================================
\section{Boundary restriction}\label{sec:boundary}
%% ===============================================================

Let $D_i\subset D$ be a smooth irreducible component,
$N_i:=\cO_{D_i}(D_i)$ its normal bundle, and $D^{(i)}:=\sum_{j\ne i}D_j$.
Write $\widehat\cO$ for the formal completion of $\cO_{\bar X}$ along
$D_i$; \emph{Zariski-locally} on $D_i$, $\widehat\cO\cong\cO_{D_i}[[t]]$
for a local equation $t$ of $D_i$ (there is in general no such
trivialisation \emph{globally} on $D_i$, precisely because $N_i$ need not
be trivial; this is the reason a coordinate-independence statement is
needed below before anything can be said globally on $D_i$).

\begin{remark}[On the status of this section]\label{rem:sec8status}
An earlier draft of \Cref{sec:boundary} proved the decomposition below
only pointwise on $D_i$ (via \Cref{thm:localstructure} applied at each
closed point) and left an explicit gap, marked at the time as
\[
   \texttt{GAP: a relative version of \Cref{thm:localstructure} over }D_i
   \texttt{ is needed to make the pointwise isomorphisms globally coherent,}
\]
to globalise it.  That gap is \textbf{not} present in the current proof
below -- it is superseded by the route described in this remark, not
silently dropped; \Cref{sec:corrections} records why the relative
statement above is hard (\Cref{lem:res0}'s Cartier-descent step does not
relativise).  The route taken here avoids the relative structure theorem
entirely, working instead directly with the diagonal.  Its logical
dependency chain is:
\[
   \Cref{prop:coordfree}\ \Longrightarrow\ \Cref{cor:globaldecomp}\
   \Longrightarrow\ \Cref{prop:blocks}\ \Longrightarrow\
   \Cref{thm:boundary},
\]
and separately \Cref{prop:jump}$\Longrightarrow$\Cref{cor:heckeexp} feeds
into \Cref{thm:boundary} as well.  \Cref{prop:jump}, \Cref{cor:heckeexp}
and \Cref{lem:normalise} are complete, elementary proofs with no gap.
\Cref{prop:coordfree}, however, is \textbf{correct in outline but not yet
formally polished} (the log PD-envelope computation needs to be written
out rigorously); consequently \Cref{cor:globaldecomp},
\Cref{prop:blocks} and \Cref{thm:boundary} all inherit that same
``proved / needs polish'' status, \emph{not} unconditional-theorem
status, until \Cref{prop:coordfree} is written up in full.  The old
relative-structure-theorem GAP is not needed any more and is not being
silently dropped: it is superseded by this route, as recorded in
\Cref{sec:corrections}.
\end{remark}

\subsection{Coordinate independence}

\begin{proposition}[\CHECK{outline correct; log PD-envelope computation
needs a formal write-up}]\label{prop:coordfree}
The operators $R_a\in\End_{\cO_{D_i}}(\bar E_0|_{D_i})$, hence the
decomposition of \Cref{def:rspec}, are independent of the choice of local
equation $t$ for $D_i$.
\end{proposition}
\begin{proof}[Proof sketch]
The $D_a$ are the dual basis of the log divided-power coordinate on the
diagonal.  Precisely, let $p_1,p_2$ be the projections of the formal
logarithmic divided-power neighbourhood of the diagonal in
$\bar X\times\bar X$, put $t_j:=p_j^*t$ and
\[
   w:=\frac{t_2}{t_1}-1 ,
\]
a well-defined element of the log-PD-envelope with divided powers
$w^{[a]}$; the logarithmic stratification is
$\epsilon(1\otimes x)=\sum_aD_a(x)\otimes w^{[a]}$.

Let $t'=ut$ with $u$ a unit.  Then
$w'=\tfrac{t'_2}{t'_1}-1=\tfrac{u_2}{u_1}(1+w)-1$.  Now $u_2-u_1$ lies in
the ideal of the diagonal, and modulo the ideal generated by $t_1$
together with the conormal directions along $D_i$ we have
$t_2-t_1=t_1w\equiv0$ and $u_2-u_1\equiv0$.  Restricting the
stratification to $D_i$ -- which is exactly the operation computing the
residues -- therefore gives $w'=w$, hence $D'_a=D_a$ on $\bar E_0|_{D_i}$.

\CHECK{this sketch needs a rigorous account of the log PD envelope and of
what ``restricting to $D_i$'' means precisely at the level of the
divided-power neighbourhood, e.g.\ a direct verification of
$D'_a-D_a\in t\,\cD(\log D)+\cD(\log D)\cdot\mathrm{Der}_{D_i}$; see
\Cref{rem:sec8status}.}
\end{proof}

\begin{corollary}\label{cor:globaldecomp}
$\bar E_0|_{D_i}=\bigoplus_{c\in\Zp}(\bar E_0|_{D_i})^c$ is a canonical
decomposition into $\cO_{D_i}$-subbundles, defined globally on $D_i$
(status as in \Cref{rem:sec8status}).
\end{corollary}
\begin{proof}
The $R_a$ are global $\cO_{D_i}$-linear endomorphisms
(\Cref{lem:leibniz,prop:coordfree}) with $R_a^p=R_a$, so the spectral
idempotents $\prod_{c'\ne c}(R_a-c')/(c-c')$ are global; multiplicities
are constant by \Cref{lem:constant}.
\end{proof}

\subsection{Residue jumps are Hecke exponents}

\begin{proposition}\label{prop:jump}
Let $M\subseteq L$ be $\cD(\log D_i)$-stable lattices in $E$ over
$\widehat\cO$ with $L/M$ supported on $D_i$.  If $L$ has $r$-spectrum
$\{c\}$ (a single value) and $M$ has $r$-spectrum $\{c'\}$, then
$a:=c'-c\in\ZZ_{\ge0}$ and $M=t^aL$.
\end{proposition}
\begin{proof}
For $x$ a $\binom{c}{\bullet}$-eigenvector, \Cref{lem:leibniz} gives
\[
   \theta_b(t^ax)=\sum_{i+j=b}D_i(t^a)\theta_j(x)
   =\sum_{i+j=b}\tbinom ai\tbinom cj\,t^ax=\tbinom{a+c}{b}t^ax
\]
by Vandermonde, so $t^aL$ has $r$-spectrum $\{c+a\}$; thus $c'-c$ is an
elementary divisor of $M\subseteq L$, in particular an integer.  If all
elementary divisors of $M\subseteq L$ equal $a$ then $t^aL\subseteq M$ and
$\mathrm{length}(L/M)=\mathrm{length}(L/t^aL)$, so $M=t^aL$.
\end{proof}

\begin{corollary}\label{cor:heckeexp}
Let $\{\bar E_n\}$ be a logarithmic lattice system and
$M_n:=\sigma_n(F^*\bar E_{n+1})$.  Then $M_n$ is $\cD(\log D_i)$-stable,
and the $r$-spectrum of $F^*\bar E_{n+1}$ is $p\cdot r(\bar E_{n+1})$.
Hence the \emph{Hecke exponents} of $\sigma_n$ are
\[
   a_n=p\,r_{n+1}-r_n\ \in\ \ZZ^N ,\qquad r_n:=r(\bar E_n).
\]
\end{corollary}
\begin{proof}
$\cD(\log)$-stability: $\partial^{[a]}(1\otimes x)=0$ unless $p\mid a$ and
$D_{pb}(1\otimes x)=t^{pb}\otimes\partial^{[b]}x=1\otimes D_b(x)$.  Thus on
$F^*\bar E_{n+1}$ the residues are $R_a=0$ for $p\nmid a$ and
$R_{pb}=R_b(\bar E_{n+1})$: the digits are shifted up by one, i.e.\
$r(F^*\bar E_{n+1})=p\,r_{n+1}$.  Now apply \Cref{prop:jump}.
\end{proof}

\begin{remark}[The Hecke exponent is \emph{not} the digit
$\dig_n(c)$]\label{rem:heckenotdigit}
An earlier draft wrote the twist below using $\dig_n(c)$ in place of
$a_n$; this is \textbf{wrong} in general and is corrected here (see
\Cref{sec:corrections} for the correction history and a worked
counterexample).  The two agree for the specific lattice system used to
\emph{define} $\cL_c$ in \Cref{sec:localstructure} (there $r_n$ is $c$
with its first $n$ digits truncated, and one checks directly that
$a_n=p\,r_{n+1}-r_n=-\dig_n(c)$ for that particular choice).  But
\Cref{lem:normalise} below normalises a general lattice system
differently -- to a \emph{single} residue value $c_n$ at each level, not
to a digit-truncation -- and for that normalisation $a_n=pc_{n+1}-c_n$
need not equal $\pm\dig_n(c)$: e.g.\ if $c_n=c$ for every $n$ (the
``constant'' normalisation) then $a_n=c(p-1)$ for all $n$, which is
independent of $n$ and generally differs from $\dig_n(c)$.  Both
normalisations give the same total twist $N_i^{\otimes c}$
(\Cref{thm:boundary}(2)); only the level-by-level bookkeeping differs.
\end{remark}

\subsection{Block decomposition of the lattices}

\begin{proposition}[status as in \Cref{rem:sec8status}]\label{prop:blocks}
Let $\bar E$ be a logarithmic lattice over $\widehat\cO$.  Then
\[
   \bar E=\bigoplus_{\bar c\in\Zp/\ZZ}\bar E\{\bar c\}
\]
canonically, where each $\bar E\{\bar c\}$ is a $\cD(\log D_i)$-stable
$\widehat\cO$-subbundle whose restriction to $D_i$ is the sum of the
components $(\bar E|_{D_i})^c$ with $c\mapsto\bar c$.  The decomposition
is preserved by every morphism of logarithmic lattices, in particular by
the $\sigma_n$.
\end{proposition}
\begin{proof}
By \Cref{cor:rigidity} the commuting $k$-linear operators $\theta_a$ on
$\bar E$ are simultaneously diagonalisable with joint characters
$\binom{c}{\bullet}$, $c\in\Zp$; write $\bar E=\bigoplus_{c}\bar E[c]$ as
$k$-vector spaces.  Each $\theta_a$ is $\cO_{D_i}$-linear
(\Cref{lem:leibniz}, using that $D_a$ involves only $\partial_t$, so
$D_b(f)=0$ for $b\ge1$ and any $f\in\cO_{D_i}$), so each $\bar E[c]$ is an
$\cO_{D_i}$-module, and by the computation in \Cref{prop:jump},
\[
   t\cdot\bar E[c]\subseteq\bar E[c+1].
\]
Therefore $\bar E\{\bar c\}:=\bigoplus_{c\mapsto\bar c}\bar E[c]$ is
$t$-stable, i.e.\ an $\widehat\cO$-submodule, and the sum over
$\bar c\in\Zp/\ZZ$ is direct and exhausts $\bar E$.  Being a direct
summand of a projective module it is projective, hence a subbundle.
Morphisms of logarithmic lattices commute with the $\theta_a$ and so
preserve the $\bar E[c]$, a fortiori the blocks.  Finally, the blocks are
determined by the restriction to $D_i$ (\Cref{cor:globaldecomp}), which
depends on \Cref{prop:coordfree} for its coordinate-freeness; hence so
does this decomposition.
\end{proof}

\begin{remark}
The finer decomposition by $c\in\Zp$ does \emph{not} lift to the lattice:
residues differing by a positive integer are resonant, and the
corresponding extension of lattices need not split.  Grouping by
$\Zp/\ZZ$ is precisely what removes resonance.  Over $K$ resonance is
invisible ($\cL_c\cong\cL_{c'}$ iff $c-c'\in\ZZ$, \Cref{rem:completemodZ}),
which is why \Cref{thm:localstructure} gives a full splitting there.
\end{remark}

\begin{lemma}[Normalisation]\label{lem:normalise}
Within a block $\bar E_n\{\bar c\}$ one may perform elementary
modifications along the subbundles $(\bar E_n|_{D_i})^{c'}$ so that the
$r$-spectrum of the block becomes a single value $c_n\in\Zp$ with
$c_n\mapsto\bar c$.  This may be done independently at each level.
\end{lemma}
\begin{proof}
An elementary modification $\ker(\bar E\to Q)$ along a subbundle $Q$ of
$\bar E|_{D_i}$ shifts the residues in $Q$ by $+1$ and leaves the others
fixed (\Cref{prop:jump}); the eigenvalues within a block differ by
integers, so finitely many such modifications equalise them.  The
$\sigma_n$ impose no constraint: for any choice of lattices,
$\sigma_n(F^*\bar E_{n+1})$ and $\bar E_n$ are commensurable lattices in
the same $K$-space.
\end{proof}

\subsection{The boundary restriction theorem}

\begin{theorem}[Boundary restriction; status as in
\Cref{rem:sec8status}]\label{thm:boundary}
Let $\cE$ be regular singular with logarithmic lattice system, normalised
as in \Cref{lem:normalise}, and fix a block $\bar c\in\Zp/\ZZ$ with
level-$n$ residue value $c_n\in\Zp$.  Put $B_n:=\bar E_n\{\bar c\}|_{D_i}$
and $a_n:=pc_{n+1}-c_n\in\ZZ$ (the Hecke exponent of
\Cref{cor:heckeexp}, \emph{not} the digit $\dig_n(c)$ -- see
\Cref{rem:heckenotdigit}). Then:
\begin{enumerate}
\item $\sigma_n$ induces isomorphisms
\[
   F_{D_i}^*\,B_{n+1}\ \isoto\ B_n\otimes N_i^{\otimes(-a_n)} ,
\]
and each $B_n$ carries a logarithmic connection along
$(D_i,D_i\cap D^{(i)})$ compatible with these.
\item Setting $e_n:=c_n$ solves the untwisting recursion
      $e_n=pe_{n+1}-a_n$; thus $(B_n)$ is an \emph{$N_i$-twisted regular
      singular stratified bundle of weight $c:=c_0$} on
      $(D_i,D_i\cap D^{(i)})$, with total twist ``$N_i^{\otimes c}$''.
\item $(B_n)$ becomes an honest regular singular stratified bundle after
      tensoring by a line bundle if and only if there exist
      $M_n\in\Pic(D_i)$ with $[M_n]=p[M_{n+1}]-a_n[N_i]$, i.e.\ iff the
      element $c\cdot[N_i]$ of the $p$-adic completion
      $\widehat{\Pic(D_i)}_p$ lies in the image of $\Pic(D_i)$.  The
      obstruction thus lies in $\widehat{\Pic(D_i)}_p/\Pic(D_i)$, and
      vanishes when $c\in\ZZ$ or $N_i$ is torsion.
\item The adjoint tractor bundle $\cA|_{D_i}$ has cancelling twists and is
      always an honest log parabolic geometry on
      $(D_i,D_i\cap D^{(i)})$.
\item $\cR_i\colon\cE\mapsto\bigl(\bar E_\bullet\{\bar c\}|_{D_i}\bigr)_{\bar c}$
      is compatible with $\otimes$, duals and pullbacks, and residues along
      $D^{(i)}$ restrict as in
      \Cref{thm:functoriality}\eqref{it:pullback}.
\end{enumerate}
\end{theorem}
\begin{proof}
(1) By \Cref{prop:blocks} the $\sigma_n$ preserve blocks; after
normalisation each block of $\bar E_n$ has a single residue value, so
\Cref{cor:heckeexp} and \Cref{prop:jump} give
$\sigma_n(F^*\bar E_{n+1}\{\bar c\})=t^{a_n}\bar E_n\{\bar c\}
=\bar E_n\{\bar c\}(-a_nD_i)$.  Restricting to $D_i$ and using
$\bigl(F^*\bar E_{n+1}\bigr)\big|_{D_i}=F_{D_i}^*\bigl(\bar E_{n+1}|_{D_i}\bigr)$
(absolute Frobenius preserves $D_i$) and
$\cO(-a_nD_i)|_{D_i}=N_i^{\otimes(-a_n)}$ gives the displayed isomorphism.
Horizontality along $D_i$: a coordinate $s$ along $D_i$ satisfies
$[\,\delta^{(m)}_t,\partial_s^{[j]}\,]=0$, so all the constructions above
are flat in the $D_i$-directions, and the logarithmic poles along
$D_j$, $j\ne i$, are preserved.

(2) $pc_{n+1}-a_n=pc_{n+1}-pc_{n+1}+c_n=c_n$.

(3) Untwisting means finding $M_n$ with $F^*(B_{n+1}\otimes M_{n+1})\cong
B_n\otimes M_n$, i.e.\ $[M_n]=p[M_{n+1}]-a_n[N_i]$.  Iterating,
$[M_0]=-\bigl(\sum_{i<n}a_ip^i\bigr)[N_i]+p^n[M_n]$, so $[M_0]$ must
reduce to $-\bigl(\sum_{i<n}a_ip^i\bigr)[N_i]$ modulo $p^n\Pic(D_i)$ for
all $n$; since $\sum_na_np^n=-c$ this says $[M_0]\mapsto c\cdot[N_i]$ in
$\widehat{\Pic(D_i)}_p$.  Conversely such an $M_0$ determines the $M_n$.

(4) and (5) are formal.
\end{proof}

\begin{example}
$\bar X=\PP^n$, $D_i$ a smooth hypersurface of degree $d$, so
$\Pic(D_i)=\ZZ$ (for $n\ge4$) and $[N_i]=d$.  \Cref{thm:boundary}(3) then
says the boundary object untwists iff $cd\in\ZZ$, i.e.\ iff
$c\in\frac1d\ZZ$ --- a strong constraint recovering \Cref{cor:extension}
when $d=1$.
\end{example}

\subsection{The dictionary with conformal infinity}\label{sec:conformal}

\begin{center}
\begin{tabular}{@{}p{0.44\textwidth}p{0.48\textwidth}@{}}
\hline
\textbf{Characteristic $0$ (Poincar\'e--Einstein)} &
\textbf{Characteristic $p$ (this paper)}\\
\hline
conformal infinity $\partial M$ & boundary component $D_i$\\
choice of defining function $\rho$ & choice of logarithmic lattice $\bar E_0$\\
conformal class (scaling of $\rho$) & indeterminacy of $r$ by $\Xst$ (\Cref{thm:welldef})\\
asymptotic expansion exponents & $p$-adic digits $R^{(m)}$ of $r$\\
conformal density of weight $w$ & $N_i^{\otimes c}$-twist, $c\in\Zp$ (\Cref{thm:boundary})\\
tractor bundle is weight $0$ & $\cA|_{D_i}$ untwisted (\Cref{thm:boundary}(4))\\
quasi-unipotent monodromy at infinity & torsion $r\in\ZZ_{(p)}$, i.e.\ tameness\\
unipotent part, weight filtration & \emph{absent locally} (\Cref{cor:nounipotent});
                                    global $\Ext$ only\\
\hline
\end{tabular}
\end{center}

\begin{remark}
The identification ``residue $=$ conformal weight'' explains the
$p$-adicity structurally: $N_i^{\otimes c}$ is meaningful inside the
$F$-divided formalism precisely for $c\in\Zp$, because the untwisting
recursion $e_n=pe_{n+1}-a_n$ (\Cref{thm:boundary}(2)) is a $p$-adic
expansion.
\end{remark}

%% ===============================================================
\section{$p$-adic middle convolution (exploratory)}\label{sec:MC}
%% ===============================================================

This section is exploratory and, apart from \Cref{con:MC,rem:middle},
does not contain unconditional theorems.  We sketch a candidate $p$-adic
middle convolution $\MC_c$, $c\in\Zp$, on $\PP^1$.  The essential input is
\Cref{thm:localstructure}: because local monodromy is semisimple, the
local invariants entering the classical theory reduce to
\emph{multiplicities} of residues rather than counts of Jordan blocks.
Whether this candidate construction actually behaves like a middle
convolution -- in particular, whether it preserves regular singularity at
all -- is the open \Cref{prob:MCrs}, on which everything else in this
section depends; see the explicit hypotheses attached to each statement
below.

\subsection{Convolution via Gauss--Manin}
Let $S=\{x_1,\dots,x_s\}\subset\AAA^1$ and $U=\AAA^1\setminus S$.  For
$c\in\Zp$ let $\cL_c$ be the rank one stratified bundle on $\GG_m$ with
residues $c$ at $0$ and $-c$ at $\infty$ (\Cref{prop:Gm}), and write
$\cL_c(y-x)$ for its pullback under $(x,y)\mapsto y-x$.

\begin{construction}\label{con:MC}
For $\cE\in\Strat^{\rs}(U)$ and $c\in\Zp$ put, for $y$ varying in
$\AAA^1\setminus S$,
\[
   \MC_c(\cE)_y:=\im\Bigl(
     H^1_{\dR,c}\bigl(\AAA^1\setminus(S\cup\{y\}),\ \cE\otimes\cL_c(y-x)\bigr)
     \longrightarrow
     H^1_{\dR}\bigl(\cdots\bigr)\Bigr).
\]
Relative de Rham cohomology of a stratified bundle along a smooth
morphism carries a Gauss--Manin stratification (Katz--Oda;
\NEEDSLITERATURECHECK{precise reference for compatibility with
$F$-divided systems}), so $\MC_c(\cE)$ is a stratified bundle on
$\AAA^1\setminus S$
\emph{wherever the cohomology sheaves are locally free of the expected
rank} -- see \Cref{prob:MClf}.
\end{construction}

\begin{remark}[Why the local invariant simplifies]\label{rem:middle}
In characteristic $0$, the definition of $j_{!*}$ requires bookkeeping of
the unipotent part of local monodromy: $\dim\cF^{I_x}$ is the
\emph{number of Jordan blocks} with eigenvalue $1$, not the multiplicity
of the eigenvalue.  By \Cref{thm:localstructure} our local monodromy is
semisimple, so, \emph{granting that the classical formula
$\dim\cF^{I_x}=\#\{j:r_{x,j}\in\ZZ\}$ makes sense at all in this
setting} (this is part of hypothesis (H-$\chi$) below),
\[
   \dim\cF^{I_x}=\#\{\,j:\ r_{x,j}\in\ZZ\,\}=:\delta_x
\]
is simply the plain multiplicity -- this is an unconditional consequence
of \Cref{thm:localstructure} \emph{given} that $\dim\cF^{I_x}$ is a
meaningful invariant here (\Cref{cor:nounipotent} already tells us there
is no unipotent part to bookkeep).  Note \Cref{cor:nosection} forbids a
canonical extension in the non-torsion range, but this is irrelevant to
\Cref{con:MC}: the cohomological definition of $\MC_c$ requires no
normalisation of $r$ into $[0,1)$.
\end{remark}

\subsection{Rank and exponents: a conditional proposition}

\begin{proposition}[\textbf{Conditional}]\label{prop:MCrank}
Assume:
\begin{description}
\item[(H-rs)] $\MC_c(\cE)$, as constructed in \Cref{con:MC}, is regular
      singular on $\AAA^1\setminus S$ for $c\notin\ZZ$ and $\cE$
      irreducible non-trivial (this is exactly \Cref{prob:MCrs});
\item[(H-$\chi$)] the positive-characteristic Euler-characteristic /
      middle-extension formula
      $\chi(\PP^1,j_{!*}\cF)=M(2-|T|)+\sum_{x\in T}\dim\cF^{I_x}$ holds
      for regular singular stratified bundles $\cF$ of rank $M$ on
      $\PP^1$ with singular locus $T$, with $j_{!*}\cF$ defined
      cohomologically as in \Cref{con:MC} and $\dim\cF^{I_x}$ as in
      \Cref{rem:middle}.
\end{description}
Then, for $c\notin\ZZ$ and $\cE$ irreducible non-trivial with $r$-spectra
$r_i$ at $x_i$ and $r_\infty$ at $\infty$, $N=\rk\cE$:
\[
   \rk\MC_c(\cE)=Ns-\sum_{i=1}^s\delta_i-\delta_\infty(c),
   \qquad
   \delta_i:=\#\{j: r_{i,j}\in\ZZ\},\quad
   \delta_\infty(c):=\#\{j: r_{\infty,j}-c\in\ZZ\}.
\]
\end{proposition}
\begin{proof}[Proof, given (H-rs) and (H-$\chi$)]
By (H-$\chi$) applied to $\cF=\cE\otimes\cL_c(y-x)$ with
$T=\{x_1,\dots,x_s,y,\infty\}$, $|T|=s+2$: the residues of $\cF$ are
$r_i$ at $x_i$, $c$ at $y$ (so $\dim\cF^{I_y}=0$ as $c\notin\ZZ$), and
$r_{\infty,j}-c$ at $\infty$.  By (H-rs), $\cF$ is regular singular, so
$j_{!*}\cF$ and $\dim\cF^{I_x}$ are meaningful, and \Cref{rem:middle}
identifies $\dim\cF^{I_x}$ with $\delta_x$.  Then
$\dim H^1_{\mathrm{mid}}=M(|T|-2)-\sum_x\dim\cF^{I_x}$ (from (H-$\chi$)
together with $\chi(\PP^1,j_!\cF)=M(2-|T|)$) gives the stated rank.
\end{proof}

\begin{remark}
Given (H-rs) and (H-$\chi$), this argument is formally complete: it uses
only \Cref{thm:localstructure} (unconditional) beyond the two hypotheses.
Neither hypothesis is established in this manuscript.  (H-rs) is
\Cref{prob:MCrs}.  (H-$\chi$) is, in effect, the assertion that the
classical Euler-characteristic bookkeeping for middle extensions of local
systems has a verbatim analogue for regular singular stratified bundles;
no proof or reference for this analogue in positive characteristic is
known to the author at the time of writing --
\NEEDSLITERATURECHECK{does such a formula appear in the literature on
$F$-divided sheaves / stratified bundles, e.g.\ in work building on
Katz--Oda?}.
\end{remark}

\subsection{Exponent rules: conjectural}

\begin{conjecture}[Exponent rules]\label{prop:MCrules}
In the situation of \Cref{prop:MCrank} (assuming (H-rs) and (H-$\chi$)
there), write $N'=\rk\MC_c(\cE)$.  The $\Zp$-valued Riemann scheme of
$\MC_c(\cE)$ is
\[
\begin{aligned}
   \text{at }x_i:&\quad \{\,r_{i,j}+c\ :\ r_{i,j}\notin\ZZ\,\}
     \ \cup\ \{\,0\ \text{with multiplicity}\ N'-(N-\delta_i)\,\},\\[2pt]
   \text{at }\infty:&\quad \{\,r_{\infty,j}-c\ :\ r_{\infty,j}-c\notin\ZZ\,\}
     \ \cup\ \{\,-c\ \text{with multiplicity}\ N'-(N-\delta_\infty(c))\,\},
\end{aligned}
\]
all equalities modulo $\ZZ$, i.e.\ up to the choice of lattice.
\end{conjecture}

\begin{remark}
Unlike \Cref{prop:MCrank}, this is \emph{not} known to follow formally
from (H-rs)+(H-$\chi$): it additionally requires the precise local
behaviour of the Gauss--Manin connection near each singular point, which
is not computed here.  What is given is Katz's classical exponent rules
\cite[Ch.~6]{Katz2}, transported formally to $\Zp$ by
$\mu=e^{2\pi i\alpha}\rightsquigarrow\alpha$ and simplified using
\Cref{rem:middle} (no Jordan-block clause).  This transport is why the
statement is recorded as a conjecture rather than a proposition: no
independent, characteristic-$p$ derivation is given.  \Cref{prob:MCintrinsic}
asks for a direct proof via the local Gauss--Manin residue; it is checked
below only on the rank-one hypergeometric example
(\Cref{ex:2F1}) and against \Cref{ex:GL2} (\Cref{ex:consistent}).
\end{remark}

\begin{example}[$p$-adic Gauss hypergeometric]\label{ex:2F1}
Let $\cE=\cL_{\alpha}(x)\otimes\cL_{\beta}(x-1)$ on
$\PP^1\setminus\{0,1,\infty\}$: rank $1$, residues $\alpha$ at $0$,
$\beta$ at $1$, $-\alpha-\beta$ at $\infty$; assume
$\alpha,\beta,\alpha+\beta+c\notin\ZZ$.  If \Cref{prop:MCrank} and
\Cref{prop:MCrules} hold in this case, then $s=2$,
$\delta_1=\delta_2=\delta_\infty(c)=0$, so
\[
   \rk\MC_c(\cE)=1\cdot2-0-0=2,
\]
with Riemann scheme
\[
   \begin{array}{c|c|c}
   0 & 1 & \infty\\\hline
   0,\ \alpha+c & 0,\ \beta+c & -c,\ -\alpha-\beta-c
   \end{array}
\]
Comparing with the classical scheme
$\bigl(\begin{smallmatrix}0,1-\gamma&0,\gamma-a-b&a,b\end{smallmatrix}\bigr)$
gives $\gamma=1-\alpha-c$, $a+b=1-\alpha-\beta-2c$, $ab$ determined; the
sum of all six exponents is $0$, consistent with
\Cref{prop:restheorem} and $\deg\bar E_0=0$ (the classical value $1$ is a
different lattice normalisation).  This is a \emph{consistency check on
the conjecture}, not a proof of it; we call the resulting object the
\emph{$p$-adic hypergeometric stratified bundle} $\cH(a,b;\gamma)$,
$a,b,\gamma\in\Zp$, understanding that its existence is conditional on
\Cref{prop:MCrank,prop:MCrules}.
\end{example}

\begin{remark}[Relation to Dwork]
A $\Zp$-parameter hypergeometric is, level by level, a compatible system
of hypergeometric-type equations with parameters
$a,b,\gamma\bmod p^{n+1}$.  This is suggestive of the shape of Dwork's
$p$-adic hypergeometric theory and of his congruences for
${}_2F_1$-coefficients \NEEDSLITERATURECHECK{compare precisely with
Dwork's Frobenius structure; not verified here}.
\end{remark}

\begin{example}[Consistency with \Cref{ex:GL2}]\label{ex:consistent}
Taking $c=\tfrac12$ (possible for $p\ne2$) and $\alpha$ suitable, the
conjectural scheme of \Cref{ex:2F1} reproduces the exponent pattern
$\{\tfrac\alpha2,\tfrac{\alpha+1}2\}$ of the pushforward along the double
cover of \Cref{ex:GL2}.  Again, this is evidence for the conjecture, not
a proof.
\end{example}

\subsection{Rigidity and the $p$-adic Katz algorithm: conjectural}

\begin{definition}
For $\cE$ regular singular on $\PP^1\setminus T$ define the
\emph{rigidity index}
\[
   \rig(\cE):=2N^2-\sum_{x\in T}\bigl(N^2-\dim Z(r_x)\bigr),
   \qquad
   \dim Z(r_x)=\sum_{\bar c\in\Zp/\ZZ}m_{x,\bar c}^2 ,
\]
where $m_{x,\bar c}$ is the multiplicity of $\bar c$ in the $r$-spectrum
at $x$.  By \Cref{thm:localstructure} the centraliser of local monodromy
really has this dimension -- no Jordan blocks intervene.  $\cE$ is
\emph{rigid} if $\rig(\cE)=2$.
\end{definition}

\begin{conjecture}\label{prop:rigMC}
$\rig(\MC_c(\cE))=\rig(\cE)$ for $c\notin\ZZ$, and $\MC_c$ preserves
irreducibility.  Moreover $\MC_{c'}\circ\MC_c=\MC_{c+c'}$ on irreducible
non-trivial objects, and $\MC_{-c}\circ\MC_c=\id$.
\end{conjecture}

\begin{remark}
This is transported from \cite[6.0]{Katz2} and depends on
\Cref{prop:MCrules} (also conjectural): granting \Cref{prop:MCrules},
rigidity-invariance is expected to follow by the same purely
combinatorial computation as in characteristic $0$, but that computation
is not carried out here, so \Cref{prop:rigMC} is recorded as a conjecture
in its own right rather than as a corollary.  A direct Gauss--Manin proof,
bypassing \Cref{prop:MCrules} entirely, would be preferable; see
\Cref{prob:MCintrinsic}.
\end{remark}

\begin{conjecture}[$p$-adic Katz algorithm]\label{conj:katzalg}
Every irreducible rigid regular singular stratified bundle on
$\PP^1\setminus T$ is obtained from a rank one object by a finite sequence
of operations $\cE\mapsto\cE\otimes\cL$ ($\cL$ of rank one) and
$\cE\mapsto\MC_c(\cE)$, $c\in\Zp$.
\end{conjecture}

\begin{remark}
The choice of $c$ in Katz's reduction step is $c=-\sum(\text{certain
residues})$.  In characteristic $0$ this requires a root of unity of the
right order in the coefficient field; here $c$ ranges over the ring $\Zp$
and such sums always exist, so the algorithm has \emph{more} room, not
less -- \emph{if} the rest of the construction goes through.  The
condition to check at each step is only $c\notin\ZZ$.
\end{remark}

\subsection{What is still missing}

\begin{problem}\label{prob:MCrs}
Does $\MC_c$ preserve regular singularity?  Equivalently: is the
Gauss--Manin stratification of a regular singular stratified bundle along
a family of affine curves again regular singular?  In characteristic $0$
this is standard; here it is open, and it is the hypothesis (H-rs) on
which \Cref{prop:MCrank} -- and hence everything else in \Cref{sec:MC} --
depends.  This is the main gap in \Cref{sec:MC}.
\end{problem}

\begin{problem}\label{prob:MClf}
Local freeness and base change for $H^1_{\dR}$ in \Cref{con:MC}: identify
hypotheses (e.g.\ on $\delta_x$) guaranteeing that $\MC_c(\cE)$ is a
vector bundle of the rank predicted by \Cref{prop:MCrank}, and that
(H-$\chi$) holds.
\end{problem}

\begin{problem}\label{prob:MCintrinsic}
Give an intrinsic proof of \Cref{prop:MCrules} by computing the local
Gauss--Manin residue, independent of the transport from
characteristic $0$, and deduce \Cref{prop:rigMC} from it directly.
\end{problem}

\begin{problem}
Combine $\MC_c$ with \Cref{thm:boundary}: is there a middle convolution in
higher dimension compatible with boundary restriction, so that the
$p$-adic Katz algorithm can be run inductively on $\dim\bar X$?  This
presupposes \Cref{prob:MCrs}.
\end{problem}

%% ===============================================================
\section{Consequences and the inductive framework}\label{sec:applications}
%% ===============================================================

\begin{corollary}[Extension criterion]\label{cor:extension}
$\cE$ extends to a stratified bundle in a neighbourhood of $D_i$ if and
only if $r_i\in\Xst$.
\end{corollary}
\begin{proof}
$(\Leftarrow)$ As in \Cref{lem:normalise}, perform elementary
modifications of $\bar E_0$ along the eigencomponents to make all
residues zero; by \Cref{cor:allzero} the $\sigma_n$ become lattice
isomorphisms, so $(\bar E_n,\sigma_n)$ is an $F$-divided sheaf near $D_i$.
$(\Rightarrow)$ Clear.
\end{proof}

\begin{corollary}\label{cor:blocks}
$\Strat^{\rs}(U)$ decomposes into blocks indexed by
$\prod_i\bigl(\Xst\otimes\Zp/\Xst\bigr)/W$; tame and non-tame blocks are
separated.
\end{corollary}

\begin{corollary}[Boundedness]
For fixed rank the possible $r$-spectra form a compact set
($\Xst\otimes\Zp$ is compact), giving quasi-compactness statements for
moduli of regular singular stratified bundles.
\TODO{make precise}
\end{corollary}

\begin{proposition}[Shape of the local-global sequence]
Local monodromy being always semisimple (\Cref{thm:localstructure}), all
non-semisimplicity of $\pi_1^{\strat,\rs}(U)$ is global.  One expects an
exact sequence
\[
   1\to \bigl(\text{``interior''}\bigr)\to
   \pi_1^{\strat,\rs}(U)\to\prod_i\cT_p ,
\]
with the tame quotient factoring through the torsion of $\Zp/\ZZ$.
\TODO{formulate and prove; this is the log analogue of Gieseker--Katz}
\end{proposition}

\subsection{Programme}
\Cref{thm:boundary} is dimension-decreasing, so induction on
$\dim\bar X$ becomes available.  Targets, in increasing order of
difficulty:
\begin{enumerate}
\item $\pi_1^{\strat,\rs}(\AAA^n)=1$ and Lefschetz-type theorems (using
      \Cref{cor:extension} and \Cref{thm:functoriality}\eqref{it:pullback}
      with $e=1$).
\item A residue theorem in $\mathrm{CH}^1(\bar X)\otimes\Zp$:
      $\sum_i\tr(r_i)[D_i]=-c_1(\bar E_0)$, and non-existence results
      derived from it.
\item A new proof of the tame Gieseker conjecture
      (\cite{EsnaultKindler}) by induction on dimension.
\item ``Twisted induction'': since $\cR_i^c$ is $N_i^{\otimes c}$-twisted,
      the inductive hypothesis must be formulated for twisted objects,
      carrying $\Pic(D_i)\otimes\Zp$ along.  Designing this correctly is
      the main structural task.
\item $p$-adic middle convolution and \Cref{prob:rigid}.
\end{enumerate}

%% ===============================================================
\section{Open problems}
%% ===============================================================

\begin{problem}[Relative local structure theorem]\label{prob:relstruct}
Prove a relative version of \Cref{thm:localstructure} over a smooth base
$S$.  This is \emph{no longer needed} to establish \Cref{thm:boundary}
(see \Cref{rem:sec8status} and \Cref{sec:corrections} for the route that
supersedes it), but is expected to matter in its own right for the
twisted induction programme of \Cref{sec:applications}.  The obstruction
to the direct approach is that \Cref{lem:res0}'s Cartier-descent step
does not relativise: the kernel $M=\ker\theta_1$ is level-$1$ horizontal
only in the $t$-direction, not in the $S$-direction.
\end{problem}

\begin{problem}[Formal write-up of coordinate independence]\label{prob:coordfreewriteup}
Write out \Cref{prop:coordfree} rigorously (the log PD-envelope
computation and the meaning of ``restricting to $D_i$'' at the level of
the divided-power neighbourhood), so that \Cref{cor:globaldecomp},
\Cref{prop:blocks} and \Cref{thm:boundary} become unconditional theorems.
This is the only remaining gap in \Cref{sec:boundary}.
\end{problem}

\begin{problem}
Since local $\Ext$ vanishes, is the global $\Ext$-object
\[
   N\in\varprojlim_n\Ext^1_{\Strat^{\rs}(U)}(\gr^W,\gr^W)
\]
the correct replacement for the unipotent part?  Is there a
``Frobenius-twisted $\mathfrak{sl}_2$'' acting on it?
\end{problem}

\begin{problem}
Compare $\cT_p=\Spec k[\Zp/\ZZ]$ with the local stratified fundamental
group scheme of $K$ including irregular objects
(cf.\ \Cref{ex:irregular}).
\end{problem}

%% ===============================================================
\section{Corrections and status of previous claims}\label{sec:corrections}
%% ===============================================================

This manuscript is a scientific record and not only a finished statement
of results; this section keeps a short account of claims made in earlier
drafts that turned out to be wrong or imprecise, why, and how they were
fixed, so that the correction history is not lost by silent editing.
Fuller working notes on these corrections are kept in
\texttt{handoff/zp-chat.md} (\S5) and in the project's daily notes; this
section is the authoritative summary inside the manuscript itself.

\subsection{Digit-wise additivity of residues (fixed prior to this
draft)}
An early version of \Cref{thm:functoriality}\eqref{it:additive} argued
tensor-additivity of the residue directly at the level of the $\Fp$-valued
digits, i.e.\ claimed
$R^{(m)}(\cE\otimes\cE')=R^{(m)}(\cE)+R^{(m)}(\cE')$.  This is false: for
example with $p=5$, $c=1$, $c'=-1$ one has $\dig_1(1)=0$ and
$\dig_1(-1)=4$, whose sum $4$ disagrees with $\dig_1(0)=0$ -- digits do
not add because of carries.  The fix, already in place throughout this
draft, is to work with the full family of operators $D_a$ (not only the
level operators $\delta^{(m)}$), use the Lucas relations
(\Cref{lem:lucas}) to reduce to a single Vandermonde identity
$\sum_{b+c=a}\binom cb\binom{c'}{c}=\binom{c+c'}a$, and only \emph{then}
extract digits via \Cref{lem:digitmap}.  The additivity is a statement
about eigenvalues in the ring $\Zp$, never about digits; see the
remark following \Cref{thm:functoriality}.

\subsection{The boundary twist exponent: digit versus Hecke exponent
(fixed in this revision)}
An earlier version of \Cref{thm:boundary} wrote the twisting isomorphism
at level $n$ as $F^*_{D_i}B_{n+1}\isoto B_n\otimes N_i^{\otimes(-\dig_n(c))}$,
using the $p$-adic digit $\dig_n(c)$ as the exponent.  This is wrong for
a general normalisation of the lattice system.  The correct exponent is
the \emph{Hecke exponent} $a_n=pc_{n+1}-c_n\in\ZZ$ of
\Cref{cor:heckeexp}, where $c_n$ is the (single, after
\Cref{lem:normalise}) residue value of the block at level $n$.  The two
agree only for the specific lattice system used to build $\cL_c$ itself
in \Cref{sec:localstructure} (where $c_n$ is $c$ with its first $n$
digits truncated, so $a_n=-\dig_n(c)$ there); for the ``constant''
normalisation $c_n\equiv c$ used in \Cref{lem:normalise}, one instead gets
the constant $a_n=c(p-1)$, generally different from $\pm\dig_n(c)$.  Both
choices telescope to the same total twist $N_i^{\otimes c}$
(\Cref{thm:boundary}(2)); only the level-by-level formula changes.  See
\Cref{rem:heckenotdigit} for the derivation and a worked comparison of
the two normalisations.

\subsection{The relative-structure GAP in \Cref{thm:boundary}: superseded,
not resolved}
An earlier draft's proof of \Cref{thm:boundary} decomposed
$\bar E_0|_{D_i}$ only pointwise on $D_i$ (applying
\Cref{thm:localstructure} at each closed point separately) and left an
explicit
\[
   \texttt{GAP: a relative version of \Cref{thm:localstructure} over }D_i
   \texttt{ is needed to make the local isomorphisms globally coherent.}
\]
An attempt to prove that relative statement directly failed: the proof of
\Cref{lem:res0} (the key step of \Cref{thm:localstructure}) uses Cartier
descent on $M=\ker\theta_1$, and $M$ is level-$1$ horizontal only in the
$t$-direction, not in the direction of the base $D_i$, so the descent does
not relativise.  This is recorded as \Cref{prob:relstruct}, which remains
open in its own right (it is expected to matter for the twisted-induction
programme of \Cref{sec:applications}) but is \textbf{no longer needed}
for \Cref{thm:boundary}.

The route actually used in the current \Cref{sec:boundary} avoids the
relative structure theorem entirely: coordinate-independence of the
residue operators $R_a$ is instead checked directly on the diagonal
(\Cref{prop:coordfree}), which combined with $\cO_{D_i}$-linearity of
$R_a$ gives the global decomposition (\Cref{cor:globaldecomp}) and the
block decomposition of lattices (\Cref{prop:blocks}) without ever
invoking a relative version of \Cref{thm:localstructure}.  This closes
the \emph{original} gap, but \Cref{prop:coordfree} itself is only correct
in outline -- the log PD-envelope computation still needs a rigorous
write-up (\Cref{prob:coordfreewriteup}) -- so \Cref{thm:boundary} should
be read as ``proved, modulo that write-up,'' not as an unconditional
theorem; see \Cref{rem:sec8status}.

\subsection{``Complete invariant'' is modulo $\ZZ$, not $\Zp$}
An earlier draft of \Cref{thm:localstructure}(2) asserted that the
$r$-spectrum itself, as an element of $\Zp$, is a complete invariant of a
regular singular stratified bundle over $K$.  This is inconsistent with
item (1) of the very same theorem: representations of the diagonalisable
group $\cT_p=\Spec k[\Zp/\ZZ]$ are graded by $\Zp/\ZZ$, not by $\Zp$.
Concretely, $\cL_c\cong\cL_{c+n}$ for every $n\in\ZZ$ (\Cref{rem:completemodZ}):
Hecke-modifying the level-$0$ lattice of the trivial bundle $\cO$ by
$n\in\Xst=\ZZ$ realises $\cO$ itself with $r$-spectrum $\{n\}$, so by the
single-eigenvalue case of \Cref{thm:localstructure}, $\cO\cong\cL_n$, and
tensoring by $\cL_c$ gives $\cL_c\cong\cL_{c+n}$.  This is the same
phenomenon, one dimension down, as $\Pic^{\strat}(\GG_m)\cong\Zp/\ZZ$ in
\Cref{prop:Gm} (already correctly stated modulo $\ZZ$ in every draft).
The corrected statement is given in \Cref{rem:completemodZ}; the
corresponding fix to the proof of \Cref{thm:localstructure} is described
in the next subsection.  This correction does not
affect the decomposition $\cE\cong\bigoplus_c\cL_c^{\oplus n_c}$, nor
\Cref{cor:extension}, both of which are about a fixed lattice, not the
isomorphism class of $\cE$ over $K$.

\subsection{The proof of \Cref{thm:localstructure}: $\Zp/\ZZ$-block
decomposition in place of eigenprojector lifting (this revision)}
Beyond the ``modulo $\ZZ$'' correction above, the proof of the
non-single-valued case of \Cref{thm:localstructure} had two defects, both
now repaired \emph{without any change to the statement of the theorem}.

\emph{(i) Eigenprojectors on the lattice.}  An earlier draft lifted the
simultaneous eigendecomposition of the $R^{(m)}$ to a
$\cD(\log)$-stable filtration of the \emph{lattice} by asserting that the
eigenprojectors ``are polynomials in the $\theta_a$, hence generate
$R$-submodules''.  This is incorrect: the $\theta_a$ are $k$-linear but
not $R$-linear (the Remark after \Cref{cor:rigidity}; indeed
$t\cdot\bar E_0^{(c)}\subseteq\bar E_0^{(c+1)}$), so a polynomial in the
$\theta_a$ is only $k$-linear and its image need not be an $R$-submodule.
The finer decomposition $\bigoplus_{c\in\Zp}E[c]$ genuinely does
\emph{not} lift to the lattice (\Cref{lem:Kblocks}; cf.\ the Remark after
\Cref{prop:blocks}).  What does lift is the coarser grouping by
$\Zp/\ZZ$-cosets: grouping the $k$-linear eigenspaces over a coset
restores stability under $t$ and $t^{-1}$ and yields a genuine
decomposition $\cE=\bigoplus_{\bar c}\cE\{\bar c\}$ of stratified bundles
over $K$ (\Cref{lem:Kblocks}, which is the specialisation to $K=k((t))$ of
the block decomposition of \Cref{prop:blocks}).  Within one block a single
residue value is reached by elementary modifications (\Cref{lem:normalise}),
reducing to the single-value case.

\emph{(ii) $\Hom=0$ versus $\Ext^1=0$.}  An earlier draft concluded that
an extension of $\cL_c$ by $\cL_{c'}$ splits from
$\Hom(\cL_c,\cL_{c'})=0$.  Vanishing of $\Hom$ does not imply vanishing of
$\Ext^1$; that step was a gap.  A first attempt to patch it treated the
resonant case $c\equiv c'\pmod\ZZ$ by claiming the $\theta_a$ are
strictly triangular on the extension, hence nilpotent, hence zero -- also
wrong, since $\theta_a$ acts on the sub-object $\cO$ as the nonzero
operator $D_a$ on $K$, and ``a self-extension has $r=0$'' is not
lattice-independent (\Cref{thm:welldef}).  The gap is now closed
structurally: once the direct-sum decomposition is established (which is
$\Ext$-free -- it uses only \Cref{lem:Kblocks} and the single-value
case), \emph{every} object is $\bigoplus_c\cL_c^{\oplus n_c}$, and with
the $\Hom$-computation the functor
$\cE\mapsto(\Hom(\cL_c,\cE))_{\bar c}$ is an equivalence of
$\Strat^{\rs}(K)$ with $\Zp/\ZZ$-graded vector spaces $=\Rep(\cT_p)$.
Semisimplicity and $\Ext^1(\cL_c,\cL_{c'})=0$ for \emph{all} $c,c'$ are
then read off the target category; concretely, an extension
$0\to\cL_{c'}\to M\to\cL_c\to0$ splits because $M\cong\bigoplus_d\cL_d^{\oplus m_d}$
and the sub-object $\cL_{c'}$ is a direct summand of the
$\bar c'$-isotypic part of $M$.

The conclusions of \Cref{thm:localstructure} --- the direct sum
decomposition, semisimplicity of $\Strat^{\rs}(K)$,
$\Ext^1(\cL_c,\cL_{c'})=0$ for all $c,c'$, and the equivalence with
$\Rep(\cT_p)$ --- are unchanged, as are the downstream statements
depending on them (\Cref{cor:nounipotent,thm:comparison,cor:blocks} and
\Cref{sec:MC}).

%% ===============================================================
\begin{thebibliography}{99}

\bibitem{DettweilerReiter} M.~Dettweiler and S.~Reiter,
  \emph{An algorithm of Katz and its application to the inverse Galois
  problem}, J. Symbolic Comput. \textbf{30} (2000), 761--798.

\bibitem{Dwork} B.~Dwork,
  \emph{$p$-adic cycles}, Publ. Math. IHES \textbf{37} (1969), 27--115.
  \TODO{also cite Dwork's hypergeometric congruences}

\bibitem{EsnaultKindler} H.~Esnault and L.~Kindler,
  \emph{Lefschetz theorems for tamely ramified coverings},
  \TODO{complete reference}

\bibitem{Gieseker} D.~Gieseker,
  \emph{Flat vector bundles and the fundamental group in non-zero
  characteristics},
  Ann. Scuola Norm. Sup. Pisa \textbf{2} (1975), 1--31.

\bibitem{Katz} N.~Katz,
  \emph{Nilpotent connections and the monodromy theorem: applications of a
  result of Turrittin},
  Publ. Math. IHES \textbf{39} (1970), 175--232.

\bibitem{Katz2} N.~Katz,
  \emph{Rigid Local Systems}, Ann. of Math. Studies \textbf{139},
  Princeton, 1996.

\bibitem{Kindler} L.~Kindler,
  \emph{Regular singular stratified bundles and tame ramification},
  \TODO{complete reference}

\bibitem{dosSantos} J.~P.~P. dos Santos,
  \emph{Fundamental group schemes for stratified sheaves},
  J. Algebra \textbf{317} (2007), 691--713.

\end{thebibliography}

\end{document}
